%% Journal of Logical and Algebraic Methods in Programming
%% Elsevier CAS single-column template
\documentclass[a4paper,fleqn]{cas-sc}

\usepackage[utf8]{inputenc}
\usepackage[english]{babel}
\usepackage[numbers,sort&compress]{natbib}
\usepackage{amsmath,amssymb,amsthm}
\usepackage{booktabs}
\usepackage{array}
\usepackage{enumitem}
\usepackage{tikz-cd}
\usepackage{algorithm}
\usepackage{algpseudocode}

%% Suppress coloured links for print-ready submission
\hypersetup{allcolors=black,colorlinks=true,hypertexnames=false}

\setlist[itemize]{topsep=2pt,itemsep=2pt}
\setlist[enumerate]{topsep=2pt,itemsep=2pt}
\sloppy

\newtheorem{definition}{Definition}
\newtheorem{proposition}{Proposition}
\newtheorem{lemma}{Lemma}
\newtheorem{theorem}{Theorem}
\newtheorem{corollary}{Corollary}
\newtheorem{remark}{Remark}

\newcommand{\Set}{\mathbf{Set}}
\newcommand{\sem}[1]{\left[\!\left[#1\right]\!\right]}
\newcommand{\D}{\mathcal{D}}
\newcommand{\Obs}{\mathsf{Obs}}
\newcommand{\Act}{\mathsf{Act}}
\newcommand{\Msg}{\mathsf{Msg}}
\newcommand{\Tool}{\mathsf{Tool}}
\newcommand{\Mem}{\mathsf{Mem}}
\newcommand{\Env}{\mathsf{Env}}
\newcommand{\Res}{\mathsf{Res}}
\newcommand{\Trace}{\mathsf{Trace}}
\newcommand{\Ev}{\mathsf{Ev}}
\newcommand{\Out}{\mathsf{Out}}
\newcommand{\cost}{\mathsf{cost}}
\newcommand{\cert}{\mathsf{cert}}
\newcommand{\RawEv}{\mathsf{RawEv}}
\newcommand{\Handoff}{\mathsf{Handoff}}
\newcommand{\Budget}{\mathsf{Budget}}
\newcommand{\Perm}{\mathsf{Perm}}
\newcommand{\Cert}{\mathsf{Cert}}
\newcommand{\Obl}{\mathsf{Obl}}
\newcommand{\Claim}{\mathsf{Claim}}
\newcommand{\Key}{\mathsf{Key}}
\newcommand{\Val}{\mathsf{Val}}
\newcommand{\Prov}{\mathsf{Prov}}
\newcommand{\Nat}{\mathbb{N}}
\newcommand{\Der}{\mathsf{Der}}
\newcommand{\Accept}{\mathsf{Accept}}
\newcommand{\TraceCheck}{\mathsf{TraceCheck}}

\title[mode=title]{A Coalgebraic and Resource-Sensitive Semantics for Tool-Augmented Agent Swarms}

\shorttitle{Coalgebraic Semantics for Tool-Augmented Agent Swarms}
\shortauthors{C.R. Ovalle}

\author[1]{Carlos Ram\'irez Ovalle}
\cormark[1]
\ead{carlosovalle@javerianacali.edu.co}
\ead[url]{https://orcid.org/0000-0002-4254-5610}
\credit{Conceptualization, Methodology, Formal analysis, Software, Writing -- original draft, Writing -- review and editing}

\affiliation[1]{organization={Departamento de Ciencias Naturales y Matem\'aticas,
                Pontificia Universidad Javeriana-Cali},
            city={Cali},
            country={Colombia}}

\cortext[1]{Corresponding author}

\begin{document}
\let\WriteBookmarks\relax
\def\floatpagepagefraction{1}
\def\textpagefraction{.001}

\begin{abstract}
Computing systems are increasingly assembled as networks of specialised, resource-governed components that use tools, maintain shared memory, delegate control, and retain execution traces. Large language model agent frameworks provide the most prominent current instance of this pattern, but the semantic status of the primitives they expose—tool permissions, handoffs, budgets, and audit traces—remains underdeveloped. This paper proposes a formal framework for tool-augmented agent swarms combining two ingredients: coalgebra, to model state-based observable dynamics, and Subexponential Linear Logic (SELL), to discipline resources, permissions, budgets, and traces. An agent is modelled as a coalgebra whose observations trigger actions, tool calls, messages, handoffs, and state updates. A swarm is a coalgebra indexed by an interaction graph and equipped with provenance-carrying shared memory and a resource ledger. The central distinction is between raw executions, which a system may emit, and certified executions, which discharge labelled obligations for graph admissibility, tool permission, budget consumption, provenance, audit evidence, and answer support. SELL subexponentials and a certification calculus provide a proof-theoretic layer distinguishing persistent, shared, local, consumable, and auditable resources. The resulting semantics supports behavioural equivalence, resource-preservation theorems for certified executions, and compositionality for separated swarms and shared append-only blackboards. The contribution is a compositional account of contemporary agentic systems as observable, resource-governed, coalgebraic processes.
\end{abstract}

\begin{highlights}
\item Coalgebra separates raw swarm behaviour from certified execution.
\item SELL labels distinguish reusable permissions from linear budgets.
\item Graph reachability determines cross-agent resource accessibility.
\item Certified runs preserve budgets, provenance, and audit evidence.
\item A 19-trace checker exercises every implemented certificate obligation.
\end{highlights}

\begin{keywords}
coalgebra \sep multi-agent systems \sep Subexponential Linear Logic \sep
resource-sensitive semantics \sep behavioural equivalence \sep
trace verification
\end{keywords}

\maketitle

\section{Introduction}

Computing systems are increasingly assembled as networks of specialised, interacting components: a coordinator delegates to specialised workers, components invoke external tools, intermediate results are written to shared memory, and execution traces are retained for auditing, debugging, and governance. Large language model (LLM) agent frameworks---such as AutoGen, LangGraph, and the OpenAI Agents SDK---provide the most prominent current instance of this pattern, with explicit primitives for multi-agent conversation, tool use, handoffs, guardrails, and tracing~\cite{wu2023autogen,li2024survey,openai2026agents,langgraph2026handoffs}.

This engineering progress raises a semantic question. What is the mathematical object denoted by a tool-augmented agent or by a swarm of such agents? Standard accounts of agents as functions from percept histories to actions,
\[
  f:\Obs^{*}\to\Act,
\]
provide a useful extensional view, but they hide internal structure: memory, tool permissions, delegation, budgets, and proof traces. Conversely, implementation frameworks expose these structures operationally, but usually without an abstract semantics capable of supporting equivalence, compositionality, and resource-sensitive verification.

This paper develops a formal framework for this gap. The guiding thesis is:
\begin{quote}
A tool-augmented agent swarm can be modelled as a coalgebraic state-based system whose transitions are governed by a resource-sensitive proof discipline.
\end{quote}
Coalgebra provides the mathematics of states, observations, transitions, and behavioural equivalence \cite{rutten2000,jacobs2016}. SELL provides a logical framework in which resources can be split into labelled zones with different structural rules, allowing one to distinguish persistent memory, shared knowledge, local agent state, tool permissions, budgets, and audit traces \cite{girard1987,nigam2009,despeyroux2017}.

The main contribution is a compositional separation between \emph{raw executions} and \emph{certified executions}. Raw executions are behaviours of a coalgebra and may include unauthorised, unaffordable, or unaudited events. Certified executions are paths whose events discharge explicit labelled obligations: graph admissibility for communication and handoffs, tool permission, exact budget consumption, provenance-carrying memory extension, fresh audit evidence, and support for final claims. This separation is intended to make tracing and guardrails semantic rather than merely operational: a trace item is not only a log entry, but evidence tied to a certified transition and to later claims.

All results specific to the proposed framework are proved explicitly in the sequel. Standard background facts about coalgebras, finality, bisimulation, and subexponential proof systems are used only as cited background; the concrete final coalgebras and preservation properties needed for this paper are derived directly.

\subsection{Contributions}

The paper makes five conceptual and technical contributions.

\begin{enumerate}
  \item It defines a coalgebraic semantics for tool-augmented agents in which observations produce events such as ordinary actions, tool calls, messages, handoffs, and final answers.
  \item It extends the single-agent account to swarms by indexing coalgebras over an interaction graph and by adding shared memory, environment state, and a resource ledger.
  \item It proposes a SELL-labelled subexponential signature and a concrete certification calculus for agentic resources, distinguishing persistent, shared, local, consumable, provenance-carrying, and trace-like assumptions, and proves that the subexponential preorder is uniquely determined by the reachability structure of the interaction graph.
  \item It states semantic properties connecting raw dynamics and certified dynamics: existence of observable behaviours under final-coalgebra assumptions, bisimulation invariance, the induced certified transition system, resource preservation for certified executions, and two compositionality results: separated swarms with disjoint resource namespaces compose without invalidating local certificates (Theorem~\ref{thm:swarm-compositionality}), and this is strictly generalised to interface composition over a shared append-only blackboard read by both components (Theorem~\ref{thm:interface-compositionality}).
  \item It provides a lightweight executable trace checker, an expanded suite of positive and negative traces, an exact checker--calculus correspondence for the implemented fragment, and a reproducible conformance scenario whose generated outputs are included as article artefacts and available in the version-controlled repository~\cite{artifact2026}.
\end{enumerate}

Taken together, these contributions apply universal coalgebra and
Subexponential Linear Logic to the correctness and resource governance of a
class of contemporary computing systems. The central formal results
are the certified safety sublanguage characterisation
(Proposition~\ref{prop:certified-safety-sublanguage}), the resource-preservation
theorem (Theorem~\ref{thm:resource-preservation}), the derivation
correspondence between the executable checker and the finite SELL calculus
(Proposition~\ref{prop:checker-adequacy}), the graph-dependent characterisation
of the subexponential preorder (Proposition~\ref{prop:graph-dependent-preorder}),
and the compositionality of certification under separated and interface swarm
composition (Theorems~\ref{thm:swarm-compositionality}
and~\ref{thm:interface-compositionality}), and the focused adequacy of the
certificate fragment (Theorem~\ref{thm:focusing-adequacy}); the prototype
operationalises and evidences these results. The coupling between the two ingredients is
structural and deliberately limited. The interaction graph that indexes the raw
coalgebra also generates the agent-to-agent part of the SELL preorder;
Proposition~\ref{prop:graph-dependent-preorder} proves that the quotient preorder
on agent labels is exactly the reachability order of the graph's strongly
connected components. Thus topology determines which local zones a focused
certificate rule may inspect. This is not a distributive law between syntax and
behaviour, and the paper does not claim such a bialgebraic coupling. The
single-pass complexity of the executable checker
(Proposition~\ref{prop:checker-complexity}) is a property of the six fixed,
locally checkable certificate schemas, not a complexity result for SELL proof
search.

The final-coalgebra constructions used below are standard in spirit, and are not presented as new theorems of universal coalgebra. They are included to make explicit which observable behaviour is being compared. The novelty is the instantiation and coupling: contemporary agent operations are represented as observable coalgebraic events, while a SELL-labelled certification calculus separates raw executions from certified executions that preserve permissions, budgets, provenance, and audit evidence. The preservation theorem is therefore a local-to-global result: once each transition rule discharges local proof obligations, arbitrary finite certified compositions inherit global resource invariants. This gives more than a post-hoc log check: it identifies the obligations that a transition must satisfy before its effects may be treated as certified evidence.

\subsection{Scope of the claim}

The claim is deliberately modest. We do not claim that all AI agents are cellular automata, finite automata, or purely symbolic systems. Nor do we claim that SELL replaces probabilistic decision models such as MDPs or POMDPs, dynamic epistemic logics, separation logics, or type-and-effect systems. Rather, we propose a semantic interface: coalgebra captures the observable dynamics of agentic execution, while SELL constrains the use and movement of resources during that execution. Probabilistic choice, reinforcement learning, epistemic update, static effect checking, and separation-style memory ownership can be added as specialised layers.

\section{Related Work}

\subsection{Coalgebra and state-based systems}

Universal coalgebra was developed as a general theory of state-based systems, including automata, streams, transition systems, and dynamical systems \cite{rutten2000,jacobs2016}. A coalgebra for an endofunctor \(F\) is a map \(c:X\to F(X)\), where \(X\) is a state space and \(F\) specifies the observable shape of one computational step. Final coalgebras provide canonical domains of observable behaviour, while bisimulations provide coinductive criteria for behavioural equivalence.

Bisimulation is also central in process semantics and model checking, where it underlies equivalence checking and refinement arguments for transition systems \cite{hennessy1985,sangiorgi2011,baier2008}. This matters for agent swarms because two different decompositions of a task may be behaviourally equivalent even when their internal prompt state, local memory, or routing structure differs.

Coalgebraic modal logic gives logical languages for specifying and distinguishing coalgebraic behaviours; Pattinson's local-consequence account and the modular framework of Cîrstea, Kurz, Pattinson, Schröder, and Venema are representative reference points \cite{pattinson2003,cirstea2011modal}. Probabilistic systems have also been studied coalgebraically. Sokolova surveys coalgebraic treatments of Markov chains, Markov processes, and probabilistic transition systems, emphasising behavioural equivalence for probabilistic dynamics \cite{sokolova2011}. Feys, Hansen, and Moss develop an algebraic and coalgebraic account of long-term values in Markov decision processes, connecting coalgebraic reasoning with optimal planning and policy improvement \cite{feys2018}.

\subsection{Coalgebraic approaches to agents and information update}

Coalgebraic ideas have already appeared in epistemic and multi-agent settings. Baltag gives a coalgebraic semantics for epistemic programs and information updates \cite{baltag2003}. Cîrstea and Sadrzadeh use coalgebraic semantics for information update in multi-agent systems without the usual change of model \cite{cirstea2007}. Carreiro, Gorín, and Schröder develop coalgebraic announcement logics, extending dynamic epistemic ideas to coalgebraic modal logic \cite{carreiro2013}. Hadzic and Chang apply coalgebra and coinduction to ontology-based multi-agent systems \cite{hadzic2008}.

These works show that coalgebraic methods are already relevant to agents, knowledge, update, and decision. The present paper differs in focus: it targets contemporary tool-augmented agent swarms, where memory, tools, handoffs, traces, and resource budgets are first-class semantic entities.

\subsection{Agentic AI frameworks and swarms}

Classical multi-agent systems already treat autonomy, social ability, reactivity, proactiveness, communication, and organisation as defining concerns \cite{wooldridge1995,shoham1993,rao1995,cohen1990,wooldridge2009}. Agent communication languages such as KQML and FIPA ACL made message structure and communicative acts explicit \cite{finin1994,fipa2002}. The present paper inherits this systems view, but its target is narrower: LLM-based agents whose execution logs expose tool calls, handoffs, shared state, and audit evidence.

The current engineering literature and documentation on LLM-based multi-agent systems emphasises role specialisation, conversation, tool use, orchestration, and delegation. AutoGen presents a multi-agent conversation framework for composing LLMs, tools, and humans in complex workflows \cite{wu2023autogen}. Recent surveys describe LLM-based multi-agent systems in terms of workflows, infrastructure, communication, profiling, cooperation, and evaluation challenges \cite{li2024survey,guo2024survey}. Tool-using language-model agents are also studied through reasoning-and-acting prompting and tool-use training regimes \cite{yao2023react,schick2023toolformer}. OpenAI's Agents SDK describes agents as LLMs equipped with instructions and tools, with handoffs, guardrails, and tracing as core primitives \cite{openai2026agents}. LangGraph similarly treats agents as graph nodes and handoffs as routing operations carrying state updates or payloads \cite{langgraph2026handoffs}.

These systems motivate the formal primitives used in this paper: agents, tools, handoffs, shared state, traces, and resource-sensitive control.

\subsection{Subexponential Linear Logic}

Linear logic is naturally resource-sensitive because assumptions are not freely reusable unless explicitly marked \cite{girard1987}. SELL refines this idea by decorating exponentials with subexponential labels. These labels are organised by a preorder and may admit different structural rules such as weakening and contraction. Nigam and Miller show that subexponentials can represent locations of multisets of formulae and increase the algorithmic expressiveness of linear logic specifications \cite{nigam2009}. Despeyroux, Olarte, and Pimentel compare Hybrid Linear Logic and SELL, showing how SELL represents modalities by labelled exponentials and an ordering among locations \cite{despeyroux2017}.

This makes SELL a natural candidate for modelling agentic resources: private memory, shared memory, tool permissions, budgets, and traces need not obey the same structural discipline.

The use of SELL here is deliberately proof-theoretic rather than merely taxonomic. Reusable permissions and trace facts are placed under structural subexponentials, while budget tokens and audit obligations remain linear. This distinction is difficult to express with a single global exponential, and it is the reason the framework uses subexponentials rather than ordinary linear logic alone.

\subsection{Process algebra, mobile calculi, and session types}

Reasoning about the behaviour of communicating and composite systems is a mature subject and the closest body of prior work to the present paper. Process calculi such as CCS and CSP model concurrent systems as terms whose communication and synchronisation are governed by structural operational rules, with bisimulation and trace or failures equivalences providing compositional notions of behavioural equality~\cite{milner1989,hoare1978csp,hennessy1985}. The \(\pi\)-calculus extends this account to \emph{mobility}: channel names are themselves communicated, so the communication topology evolves dynamically~\cite{milner1992picalculus,sangiorgi2011}. These calculi already supply the conceptual vocabulary that the present paper reuses for swarms---communication between components, delegation of capabilities, and behavioural equivalence as the criterion of sameness---and the interaction graph, the message and handoff events, and the swarm bisimulation used below are, in this respect, deliberately classical rather than novel.

The deepest link between linear logic and concurrency is the \emph{propositions-as-sessions} correspondence. Caires and Pfenning interpret intuitionistic linear-logic propositions as session types, proofs as \(\pi\)-calculus processes, and cut elimination as communication; under the assumptions of that calculus, well-typed processes enjoy protocol fidelity and global progress properties~\cite{caires2010sessions}. Wadler's reformulation relates classical linear logic directly to a session-typed calculus~\cite{wadler2014sessions}, both building on the session-type discipline of Honda and collaborators~\cite{honda1998,huttel2016}. This established line shows that linear logic is already a recognised semantic foundation for disciplined communication, which makes an explicit comparison indispensable.

The comparison also locates the present contribution precisely. In the propositions-as-sessions systems above, a linear proposition types a channel together with its dual, so the logical discipline governs the interaction protocol. The framework developed here uses linear logic for a different purpose. It does not type channels, require duality, or establish deadlock-freedom; messages and handoffs are constrained only by graph admissibility, which is strictly weaker than session fidelity. Instead, the \emph{subexponential} layer stratifies resources into labelled zones with different structural status---reusable tool permissions and trace evidence under structural subexponentials, consumable budgets and audit obligations kept linear---and certification ties these resources to quantitative budget accounting, append-only provenance memory, and answer support. Basic session types abstract from these concerns unless enriched with quantitative or refinement information. Conversely, the present calculus says nothing about two endpoints following dual protocols. The two disciplines are therefore complementary: session types govern the \emph{shape of interaction} along a channel, whereas the present subexponential calculus governs the \emph{consumption, reuse, and provenance of resources} along a certified execution.

Resource-aware session types narrow this gap by tracking work or time inside the session-typed proposition~\cite{das2018}. They share the present concern with quantitative accounting, but they fold cost into the behavioural type of a channel, whereas here the budget, permission, provenance, and audit ledger is a separate component of the coalgebraic state, decoupled from the behavioural functor and from any notion of channel duality. This decoupling is precisely what lets raw coalgebraic behaviour and certified behaviour be separated, with the certified executions forming a prefix-closed safety sublanguage of the raw behaviour (Propositions~\ref{prop:certified-subcoalgebra} and~\ref{prop:certified-safety-sublanguage}) rather than a typing discipline on process terms.

Finally, the abstract relationship between operational behaviour and algebraic structure is captured by the bialgebraic semantics of Turi and Plotkin, where a distributive law of syntax over behaviour yields well-behaved operational specifications and bisimulation congruences~\cite{turi1997}. The present paper deliberately does not give a bialgebraic or GSOS-style specification of swarm syntax: the SELL certification layer is not a distributive law generating the transition structure, but a resource-sensitive admissibility discipline that carves a certified subcoalgebra out of an already-given raw coalgebra. Coupling SELL-indexed resources to a distributive law, so that certification becomes a congruence for arbitrary interface composition, is a natural next step; the composition results here (Theorems~\ref{thm:swarm-compositionality} and~\ref{thm:interface-compositionality}) cover separated swarms and a shared append-only blackboard with cross-component reads, but not a bialgebraic congruence for shared mutable state.

Table~\ref{tab:formal-comparison} summarises the comparison along the axes that
matter for the editor's concern. The table does not claim that one formalism
subsumes another; it records different verification targets.

\begin{table}[htbp]
\centering
\scriptsize
\begin{tabular}{>{\raggedright\arraybackslash}p{0.17\linewidth}>{\raggedright\arraybackslash}p{0.24\linewidth}>{\raggedright\arraybackslash}p{0.24\linewidth}>{\raggedright\arraybackslash}p{0.21\linewidth}}
\toprule
Formalism & Primary object & Typical guarantee & Relation to this work \\
\midrule
CCS, CSP, \(\pi\)-calculus & communicating process and operational transition system & behavioural equivalence, traces, mobility, or synchronisation properties & supplies the classical communication and bisimulation vocabulary \\
Session types & channel protocol with dual endpoints & protocol fidelity and progress properties under the chosen type system & stronger on protocols; does not by itself provide this ledger discipline \\
Resource-aware session types & protocol enriched with potential or cost & protocol fidelity together with quantitative bounds & closest quantitative relative; cost remains integrated with the channel type \\
Present framework & raw coalgebra plus SELL-indexed certificate ledger & finite-run safety for permissions, budgets, provenance, audit evidence, and answer support & weaker on protocols; separates behaviour from resource certification \\
Bialgebraic semantics & syntax and behaviour connected by a distributive law & congruence of behavioural equivalence & stronger syntax--behaviour coupling; not provided here \\
\bottomrule
\end{tabular}
\caption{Comparison with formalisms for composite communicating systems.}
\label{tab:formal-comparison}
\end{table}

\subsection{Proof-carrying evidence and executable checkers}

The executable part of this paper is closer to a proof checker or certifying algorithm than to an empirical benchmark. Certifying algorithms produce outputs together with witnesses that can be checked independently \cite{mcconnell2011certifying}; proof-carrying code accepts code only when accompanied by a checkable safety proof \cite{necula1997pcc}. Certified proof checkers for SAT and related proof formats similarly distinguish generated evidence from the smaller algorithm that validates it \cite{cruzfilipe2017resolution}. Fully verified systems such as CompCert go further by mechanising large correctness proofs in a proof assistant \cite{leroy2009compcert}. The present artefact is intentionally at a lighter level: it specifies an abstract finite trace-checking function \(\TraceCheck\), proves exact correspondence with the six operational certificate rules, and provides \texttt{prototype/checker.py} as a reproducible reference implementation. It does not claim that the Python implementation itself is mechanically verified.

\subsection{Alternative formal tools}

Several neighbouring formalisms solve different parts of the problem. Dynamic epistemic logic is a natural language for announcements, belief change, and information update \cite{vanditmarsch2007}; it is less directly aimed at linear budgets or tool permissions. Separation logic gives precise reasoning principles for mutable heap ownership and frame conditions \cite{reynolds2002}, but its native objects are heaplets rather than multi-agent handoff traces. Type-and-effect systems can statically approximate which tools or effects a program may use \cite{lucassen1988,nielson1999}, and algebraic effects give a semantic account of operations and handlers \cite{plotkin2003}, but these approaches usually abstract away from run-time provenance and audit obligations. Provenance models such as W3C PROV and database provenance give a vocabulary for source history and derivation \cite{moreau2013prov,buneman2001,cheney2009}; the present framework uses that intuition inside a resource-sensitive transition discipline rather than as a standalone metadata model. MDPs and POMDPs model optimal action under uncertainty \cite{puterman1994,kaelbling1998}; they are essential for policy and planning questions, but do not by themselves provide proof objects for resource use. Guardrails and tracing in engineering frameworks are operational checks and records; the aim here is to connect such checks to compositional obligations whose preservation follows across certified executions. The present framework should therefore be read as complementary to these approaches: coalgebra supplies observable dynamics, while SELL supplies labelled resource control over executions.

\section{Preliminaries}

\subsection{Agents as external behaviours}

Let \(\Obs\) be a set of observations and \(\Act\) a set of actions. A classical extensional representation of an agent is a function
\[
  f:\Obs^{*}\to\Act,
\]
where \(f(w)\) is the action selected after the observation history \(w\). This view is maximally abstract, but it hides the internal mechanisms by which the action is produced.

A state-based representation introduces an internal state space \(X\), an output function, and a transition function. In deterministic Moore style one writes
\[
  o:X\to\Act,
  \qquad
  t:X\times\Obs\to X.
\]
In deterministic Mealy style one writes
\[
  m:X\times\Obs\to \Act\times X.
\]
Both are coalgebraic.

\subsection{Coalgebras}

\begin{definition}[Coalgebra]
Let \(F:\Set\to\Set\) be an endofunctor. An \(F\)-coalgebra is a pair \((X,c)\), where \(X\) is a set of states and
\[
  c:X\to F(X)
\]
is a transition-observation map.
\end{definition}

For the Moore functor \(F(X)=\Act\times X^{\Obs}\), a coalgebra is a map
\[
  c:X\to\Act\times X^{\Obs}.
\]
It assigns to each state an observable action and a continuation for each possible observation.

For the Mealy functor \(G(X)=(\Act\times X)^{\Obs}\), a coalgebra is a map
\[
  m:X\to(\Act\times X)^{\Obs}.
\]
It assigns, for each state and observation, an action and a next state.

\subsection{Final behaviour}

\begin{proposition}[Final Moore semantics]\label{prop:final-moore-semantics}
Let \(F(X)=\Act\times X^{\Obs}\). The set
\[
  B=\Act^{\Obs^{*}}
\]
with coalgebra structure \(\zeta:B\to\Act\times B^{\Obs}\) defined by
\[
  \zeta(b)=\bigl(b(\varepsilon),\,p\mapsto b_p\bigr),
  \qquad
  b_p(w)=b(pw),
\]
is a final \(F\)-coalgebra. Hence every \(F\)-coalgebra \(c:X\to\Act\times X^{\Obs}\) induces a unique behaviour map
\[
  \sem{-}:X\to\Act^{\Obs^{*}}.
\]
\end{proposition}

\begin{proof}
The set \(B=\Act^{\Obs^{*}}\) with the structure \(\zeta\) is the standard final coalgebra of the deterministic Moore functor \(F(X)=\Act\times X^{\Obs}\); the construction of the behaviour map and its uniqueness are part of the standard theory of coalgebras and Moore automata, and we refer to it rather than reproving it~\cite{rutten2000,jacobs2016}. The induced map \(\sem{-}\) is the unique function satisfying \(\sem{x}(\varepsilon)=o(x)\) and \(\sem{x}(pw)=\sem{t_x(p)}(w)\), where \(c(x)=(o(x),t_x)\).
\end{proof}

This map sends an internal state to its complete observable behaviour.

\subsection{Subexponential linear logic}\label{subsec:sell-prelim}

The certification layer of this paper is expressed in subexponential linear logic (SELL), a refinement of linear logic in which the exponential modality \(!\) is replaced by a family of labelled modalities \(!^{\ell}\) that may obey different structural rules~\cite{girard1987,nigam2009,despeyroux2017}. We recall the ingredients used in the certification calculus.

\begin{definition}[Subexponential signature]\label{def:subexp-signature}
A \emph{subexponential signature} is a triple \(\Sigma=(I,\preceq,U)\), where \(I\) is a set of \emph{subexponential labels}, \(\preceq\) is a preorder on \(I\), and \(U\subseteq I\) is an upward-closed set of labels admitting the structural rules of weakening and contraction: if \(u\in U\) and \(u\preceq v\), then \(v\in U\). A label \(\ell\in U\) is \emph{unbounded}: assumptions stored at \(\ell\) may be copied or discarded freely. A label \(\ell\notin U\) is \emph{linear}: its assumptions may be neither duplicated nor weakened. The preorder \(\preceq\) controls promotion, that is, which labelled assumptions a rule focused at a given label may use.
\end{definition}

\begin{definition}[Linear and subexponential contexts]\label{def:contexts}
A SELL judgement separates assumptions into two contexts. The \emph{linear context} \(\Delta\) is a finite multiset of formulae subject to no structural rule, so that each formula of \(\Delta\) must be used exactly once. The \emph{subexponential context} \(\Theta\) is a finite multiset of \emph{labelled} formulae, written \(F@\ell\) or \(!^{\ell}F\) with \(\ell\in I\); a labelled formula may be contracted or weakened precisely when its label lies in \(U\), and \(\preceq\) determines when it may be accessed by a rule focused at another label. We write \(\Theta_{\ell}\) for the part of \(\Theta\) carrying label \(\ell\), and a judgement has the form \(\Delta;\Theta\vdash G\).
\end{definition}

Thus \(\Delta\) carries consumable resources, while \(\Theta\) carries resources whose reusability is decided label-by-label by \((\preceq,U)\). The full proof system of SELL---with its labelled promotion, dereliction, and cut rules---is standard~\cite{nigam2009,despeyroux2017} and is not needed in detail here: the certification calculus of Section~\ref{subsec:sell-sequents-certificates} uses only a fixed family of rule schemas over a concrete signature \(\Sigma_{\Gamma}\), whose labels, preorder, and unbounded set are defined in Section~\ref{sec:sell-layer}.


\section{Tool-Augmented Agents as Coalgebras}

The preceding account must be refined for modern agents, since their observable outputs are not merely ordinary actions. They may call tools, send messages, delegate control, emit a final answer, update memory, or emit audit evidence. We therefore distinguish between \emph{raw events}, which may occur in an implementation trace, and \emph{well-typed events}, which are compatible with the swarm interface.

\begin{definition}[Local event interface]
For an agent \(i\), let
\[
  \Ev_i
  =
  \Act_i
  + \mathsf{ToolCall}_i
  + \Msg_i
  + \Handoff_i
  + \mathsf{Answer}_i
  + \mathsf{TraceItem}_i.
\]
The coproduct tags the observable kind of event. Thus a tool call is not identified with a tool name alone: it is an invocation carrying at least a tool identifier and a payload of arguments.
\end{definition}

\begin{definition}[Tool-augmented agent]
A deterministic tool-augmented agent is a Mealy-style coalgebra
\[
  c_i:X_i\to(\Ev_i\times X_i)^{\Obs_i}.
\]
Equivalently, for \(x\in X_i\) and \(o\in\Obs_i\), one has \(c_i(x)(o)=(e,x')\), where \(e\in\Ev_i\) is the emitted event and \(x'\in X_i\) is the next internal state.
\end{definition}

The state space \(X_i\) may contain prompt state, scratchpad, local memory, retrieved context, pending goals, tool outputs, and other implementation-level data. Coalgebra abstracts from these implementation details while preserving the observable event-and-continuation structure.

\begin{remark}[Probabilistic agents]
A stochastic version is obtained by replacing deterministic continuations with a distribution functor:
\[
  c_i:X_i\to \D(\Ev_i\times X_i)^{\Obs_i},
\]
where \(\D\) may be the finite-support distribution monad. This gives a coalgebraic interface to probabilistic transition systems, MDP-like dynamics, and stochastic policies. The present paper keeps the deterministic presentation because the resource theorem concerns certified finite executions.
\end{remark}

\subsection{Handoffs}

A handoff transfers control, context, or responsibility from one agent to another. If \(V\) is the set of agents, a handoff from \(i\) to \(j\) carrying payload \(p\) is written
\[
  \mathsf{handoff}_{ij}(p)\in\Handoff_i.
\]
At the local coalgebraic level this is just an event. At the swarm level it is admissible only if the interaction topology and the resource discipline permit it.

\section{Agent Swarms}

An agent swarm is not merely a product of isolated agents. Its essential structure is interaction: agents communicate, share memory, delegate control, compete for tools, and consume global resources. Formally, a swarm is a family of systems parameterised by the agent set and interaction graph. The three-agent example used later is only one finite instantiation; the definitions below are stated for an arbitrary set of agents~\(V\).

\begin{definition}[Interaction graph]
An interaction graph is a directed graph
\[
  \Gamma=(V,E_{\Gamma}),
\]
where \(V\) is the set of agents and \((i,j)\in E_{\Gamma}\) means that agent \(i\) may send a message or hand off control to agent \(j\).
\end{definition}

For a fixed graph \(\Gamma\), define the raw event interface
\[
\begin{aligned}
  \RawEv_{\Gamma}={}&
  \sum_{i\in V}\Act_i
  +\sum_{i\in V}\mathsf{ToolCall}_i
  +\sum_{i,j\in V}(\Msg_{ij}+\Handoff_{ij}) \\
  &+\mathsf{MemUpdate}+\mathsf{EnvUpdate}+\mathsf{Answer}+\mathsf{TraceItem}.
\end{aligned}
\]
The well-typed graph-compatible events form a subset \(\Ev_{\Gamma}\subseteq\RawEv_{\Gamma}\) obtained by retaining only those message and handoff events whose endpoints lie in \(E_{\Gamma}\). Tool permissions and budgets are not built into this type: they are checked by the resource layer.

Shared memory is treated as a provenance-carrying finite partial map
\[
  M:\Key\rightharpoonup \Val\times\Prov.
\]
If \(M(k)=(v,p)\), then key \(k\) stores value \(v\) together with provenance record \(p\), for instance the tool call, source identifier, and trace item by which the value entered the shared context. We write \(M\sqsubseteq M'\) when \(M'\) extends \(M\) without changing any previously recorded binding. The deterministic core uses this append-only discipline for shared source keys; mutable scratchpads and overwritten local state remain inside the local state spaces \(X_i\).

Let the global state space be
\[
  X_{\Gamma}=Y_{\Gamma}\times\Res,
  \qquad
  Y_{\Gamma}=\prod_{i\in V}X_i\times \Mem\times\Env.
\]
The component \(Y_{\Gamma}\) contains local agent states, shared memory, and environment state; the \(\Res\) component stores the resource ledger. We use a ledger of the form
\[
  R=(\beta,\pi,\tau,\kappa),
\]
where \(\beta:V\to\Nat\) gives the available linear budget per agent, \(\pi\subseteq V\times\Tool\) records reusable tool permissions, \(\tau\subseteq\Trace\) is the emitted audit log, and \(\kappa\subseteq\Claim\times\Key\) records certified claims together with the shared-memory key that supports each claim. Thus certification is not just a Boolean flag on a sentence: it binds a claim to a source in the shared memory.

\begin{definition}[Raw swarm coalgebra]
Given an interaction graph \(\Gamma\), a deterministic raw swarm coalgebra is a map
\[
  C:X_{\Gamma}\to(\RawEv_{\Gamma}\times X_{\Gamma})^{\Obs_{\Gamma}}.
\]
For \(x\in X_{\Gamma}\) and \(o\in\Obs_{\Gamma}\), we write
\[
  x\xrightarrow[o]{e}_{C}x'
\]
when \(C(x)(o)=(e,x')\).
\end{definition}

The coalgebra describes the events an implementation may emit. It deliberately does not guarantee that every emitted event is permitted, affordable, or auditable. Those properties are captured by certification.

\begin{definition}[Certified transition]
Let \(\Sigma_{\Gamma}\) be a SELL subexponential signature for \(\Gamma\) in the sense of Definition~\ref{def:subexp-signature}, with concrete labels and preorder fixed in Section~\ref{sec:sell-layer}. A certification relation is a judgement
\[
  \Delta;\Theta\vdash_{\Sigma_{\Gamma}} e:(y,R)\leadsto(y',R'),
\]
where \(\Delta\) is the linear context and \(\Theta\) the subexponential context in the sense of Definition~\ref{def:contexts}, \(y,y'\in Y_{\Gamma}\), \(R,R'\in\Res\), and \(e\in\RawEv_{\Gamma}\). Write \(\mathsf{AdmCtx}_{\Gamma}((y,R),\Delta,\Theta)\) when these contexts are exactly the ledger interpretation fixed in Section~\ref{subsec:sell-sequents-certificates}, together with the reusable certificate program, relevant policy zones, and the linear obligation for the event. A raw transition \((y,R)\xrightarrow[o]{e}_{C}(y',R')\) is certified only if its non-ledger and ledger components satisfy the judgement with such an admissible source context; assumptions cannot be chosen freely to manufacture a certificate.
\end{definition}

Thus the semantic object of interest is not just \(C\), but the pair
\[
  (C,\vdash_{\Sigma_{\Gamma}}),
\]
which separates possible dynamics from certified dynamics.

It is useful to name the certified subdynamics. For a raw swarm coalgebra \(C\), define the certified transition relation
\[
  (y,R)\overset{e,o}{\Longrightarrow}_{\cert}(y',R')
\]
iff \(C(y,R)(o)=(e,(y',R'))\) and there is a certificate judgement \(\Delta;\Theta\vdash_{\Sigma_{\Gamma}}e:(y,R)\leadsto(y',R')\). Thus certification turns the total raw coalgebra into a labelled transition system restricted to admissible certificate steps. If every raw transition is certified, this restricted system is again total; Proposition~\ref{prop:certified-subcoalgebra} below proves that certifiable states form an \(S_{\Gamma}\)-subcoalgebra when closed under certified successors, and that the inclusion into the raw system is a coalgebra morphism. The subcoalgebra need not be proper: it may coincide with the raw coalgebra or be empty. In the more realistic case, where some raw implementation steps are rejected, the certified behaviour is a partial LTS whose finite paths are precisely the executions used in the preservation theorem.

\begin{proposition}[Final semantics for deterministic swarms]\label{prop:final-swarm-semantics}
Let
\[
  S_{\Gamma}(X)=(\RawEv_{\Gamma}\times X)^{\Obs_{\Gamma}}.
\]
In \(\Set\), this functor has a final coalgebra isomorphic to
\[
  \RawEv_{\Gamma}^{\Obs_{\Gamma}^{+}},
\]
where \(\Obs_{\Gamma}^{+}\) is the set of non-empty finite observation words. Hence every raw swarm coalgebra \(C:X_{\Gamma}\to S_{\Gamma}(X_{\Gamma})\) induces a unique behaviour map
\[
  \sem{-}_{\Gamma}:X_{\Gamma}\to\RawEv_{\Gamma}^{\Obs_{\Gamma}^{+}}.
\]
\end{proposition}

\begin{proof}
The functor \(S_{\Gamma}(X)=(\RawEv_{\Gamma}\times X)^{\Obs_{\Gamma}}\) is the deterministic Mealy functor, whose final coalgebra is standard; the carrier \(\RawEv_{\Gamma}^{\Obs_{\Gamma}^{+}}\) and the uniqueness of the induced behaviour map are instances of the coalgebraic theory of Mealy machines, which we cite rather than reprove~\cite{rutten2000,bonsangue2008,jacobs2016}. Concretely, \(\sem{-}_{\Gamma}\) is the unique function such that, whenever \(C(x)(o)=(e,x')\),
\[
  \sem{x}_{\Gamma}(o)=e
  \qquad\text{and}\qquad
  \sem{x}_{\Gamma}(ow)=\sem{x'}_{\Gamma}(w)
  \quad(w\in\Obs_{\Gamma}^{+}).
\]
These characteristic equations, which say that \(\sem{x}_{\Gamma}(u)\) is the last event emitted after reading \(u\), are the only properties of the behaviour map used in the sequel.
\end{proof}

Although \(\sem{x}_{\Gamma}(u)\) returns the last event emitted after the observation word \(u\), it does not discard trace information. Given a concrete input word \(u=o_1\cdots o_n\), the finite observable trace is recovered from the behaviour map by evaluating it on prefixes:
\[
  \bigl(\sem{x}_{\Gamma}(o_1),
  \sem{x}_{\Gamma}(o_1o_2),\ldots,
  \sem{x}_{\Gamma}(o_1\cdots o_n)\bigr).
\]
The audit component \(\tau\) in the ledger records which of these emitted events carry certified trace identifiers. Thus coalgebraic behaviour supplies the observable event stream, while the certification layer determines which stream prefixes are auditable and resource-safe.

\subsection{Bisimulation}

\begin{definition}[Swarm bisimulation]
Let \((X,C)\) and \((Y,D)\) be two raw swarm coalgebras for the same functor \(S_{\Gamma}\). A relation \(R\subseteq X\times Y\) is a bisimulation if whenever \((x,y)\in R\), then for every observation \(o\in\Obs_{\Gamma}\), the transitions
\[
  C(x)(o)=(e,x'),
  \qquad
  D(y)(o)=(e',y')
\]
satisfy \(e=e'\) and \((x',y')\in R\).
\end{definition}

\begin{proposition}[Behavioural invariance]\label{prop:behavioural-invariance}
If \((x,y)\in R\) for some swarm bisimulation \(R\), then
\[
  \sem{x}_{\Gamma}=\sem{y}_{\Gamma}.
\]
\end{proposition}

\begin{proof}
Let \(u\in\Obs_{\Gamma}^{+}\). We prove by induction on the length of \(u\) that
\[
  \sem{x}_{\Gamma}(u)=\sem{y}_{\Gamma}(u)
\]
whenever \((x,y)\in R\). If \(u=o\) has length one, write
\[
  C(x)(o)=(e,x'),
  \qquad
  D(y)(o)=(e',y').
\]
Since \(R\) is a bisimulation, \(e=e'\). By the definition of the final behaviour map in Proposition~\ref{prop:final-swarm-semantics},
\[
  \sem{x}_{\Gamma}(o)=e=e'=\sem{y}_{\Gamma}(o).
\]
Now suppose \(u=ow\), where \(w\in\Obs_{\Gamma}^{+}\). Again write
\[
  C(x)(o)=(e,x'),
  \qquad
  D(y)(o)=(e',y').
\]
Bisimulation gives \(e=e'\) and \((x',y')\in R\). The final behaviour map satisfies
\[
  \sem{x}_{\Gamma}(ow)=\sem{x'}_{\Gamma}(w),
  \qquad
  \sem{y}_{\Gamma}(ow)=\sem{y'}_{\Gamma}(w).
\]
By the induction hypothesis applied to \((x',y')\in R\) and the shorter word \(w\),
\[
  \sem{x'}_{\Gamma}(w)=\sem{y'}_{\Gamma}(w).
\]
Therefore \(\sem{x}_{\Gamma}(ow)=\sem{y}_{\Gamma}(ow)\). Since this holds for every non-empty observation word, the two functions \(\sem{x}_{\Gamma}\) and \(\sem{y}_{\Gamma}\) are equal.
\end{proof}

\subsection{Certified Subcoalgebra and Certified Bisimulation}

The two notions of certified dynamics and bisimulation can now be made precise at the coalgebraic level.

\begin{proposition}[Certified states form a subcoalgebra]\label{prop:certified-subcoalgebra}
Let \(C:X_{\Gamma}\to S_{\Gamma}(X_{\Gamma})\) be a raw swarm coalgebra. Call a state \(x\in X_{\Gamma}\) \emph{certifiable} if every raw transition from \(x\) admits a certificate derivation:
\[
  X_{\cert}=\bigl\{x\in X_{\Gamma}\mid\forall o\in\Obs_{\Gamma},\;
  C(x)(o)=(e,x')
  \Rightarrow
  \exists\,\Delta;\Theta.\;
  \mathsf{AdmCtx}_{\Gamma}(x,\Delta,\Theta)\ \land\
  \Delta;\Theta\vdash_{\Sigma_{\Gamma}}e:x\leadsto x'
  \bigr\}.
\]
If \(X_{\cert}\) is closed under certified successor states, that is, whenever \(x\in X_{\cert}\) and \(x\overset{e,o}{\Longrightarrow}_{\cert}x'\) then \(x'\in X_{\cert}\), then:
\begin{enumerate}
  \item The restriction \(C_{\cert}:=C|_{X_{\cert}}:X_{\cert}\to S_{\Gamma}(X_{\cert})\) is a well-defined \(S_{\Gamma}\)-coalgebra.
  \item The inclusion \(\iota:X_{\cert}\hookrightarrow X_{\Gamma}\) is a coalgebra morphism from \((X_{\cert},C_{\cert})\) to \((X_{\Gamma},C)\).
  \item For every \(x\in X_{\cert}\), the unique behaviour maps satisfy \(\sem{x}_{C_{\cert}}=\sem{\iota(x)}_{\Gamma}\); certified states have identical observable behaviours whether viewed inside the subcoalgebra or inside the raw system.
\end{enumerate}
\end{proposition}

\begin{proof}
For~(1), let \(x\in X_{\cert}\) and \(o\in\Obs_{\Gamma}\), and write \(C(x)(o)=(e,x')\). By definition of \(X_{\cert}\) the transition has a certificate derivation, and by the closure hypothesis \(x'\in X_{\cert}\). Hence \((e,x')\in\RawEv_{\Gamma}\times X_{\cert}\), so the map \(C_{\cert}(x)(o):=(e,x')\) lands in \(S_{\Gamma}(X_{\cert})\).

For~(2), let \(x\in X_{\cert}\) and \(o\in\Obs_{\Gamma}\), and write \(C_{\cert}(x)(o)=(e,x')\) with \(x'\in X_{\cert}\). The functor \(S_{\Gamma}\) acts on morphisms by post-composition on the state component:
\[
  \bigl(S_{\Gamma}(\iota)\circ C_{\cert}\bigr)(x)(o)
  =(e,\iota(x'))
  =(e,x')
  =C(x)(o)
  =(C\circ\iota)(x)(o).
\]
Since this holds for every \(x\in X_{\cert}\) and \(o\in\Obs_{\Gamma}\), the diagram \(S_{\Gamma}(\iota)\circ C_{\cert}=C\circ\iota\) commutes, so \(\iota\) is a coalgebra morphism.

For~(3), by Proposition~\ref{prop:final-swarm-semantics}, the unique behaviour map \(\sem{-}_{\Gamma}\) is the unique morphism to the final \(S_{\Gamma}\)-coalgebra \((B_{\Gamma},\zeta)\). Since \(\iota\) is a coalgebra morphism from \((X_{\cert},C_{\cert})\) to \((X_{\Gamma},C)\), the composition \(\sem{-}_{\Gamma}\circ\iota:X_{\cert}\to B_{\Gamma}\) is a coalgebra morphism from \((X_{\cert},C_{\cert})\) to \((B_{\Gamma},\zeta)\). By finality this equals \(\sem{-}_{C_{\cert}}\), giving \(\sem{x}_{C_{\cert}}=\sem{\iota(x)}_{\Gamma}\) for every \(x\in X_{\cert}\).
\end{proof}

Proposition~\ref{prop:certified-subcoalgebra} makes the informal remark earlier in this section precise: certification is not merely an external filter on an existing LTS, but carves out a genuine coalgebraic substructure embedded by a morphism. Two different resource policies for the same raw system yield two different subcoalgebras, whose inclusion morphisms into the raw system allow their observable behaviours to be compared via the final coalgebra.

The notion of bisimulation extends naturally to the total certified dynamics on
\(X_{\cert}\). Restricting to certifiable states is essential: on the partial LTS
outside \(X_{\cert}\), a one-sided simulation could hold vacuously at a state
with no certified successor and would not imply equality of raw behaviour.

\begin{definition}[Certified bisimulation]\label{def:certified-bisimulation}
Assume the closure condition of Proposition~\ref{prop:certified-subcoalgebra}.
A relation \(R\subseteq X_{\cert}\times X_{\cert}\) is a
\emph{certified bisimulation} if, whenever \((x,y)\in R\), both of the
following conditions hold for every \(o\in\Obs_{\Gamma}\):
\begin{enumerate}
  \item if \(x\overset{e,o}{\Longrightarrow}_{\cert}x'\), then there is
  \(y'\in X_{\cert}\) with
  \(y\overset{e,o}{\Longrightarrow}_{\cert}y'\) and \((x',y')\in R\);
  \item if \(y\overset{e,o}{\Longrightarrow}_{\cert}y'\), then there is
  \(x'\in X_{\cert}\) with
  \(x\overset{e,o}{\Longrightarrow}_{\cert}x'\) and \((x',y')\in R\).
\end{enumerate}
Two states \(x,y\in X_{\cert}\) are \emph{certified bisimilar}, written
\(x\sim_{\cert}y\), if some certified bisimulation contains \((x,y)\).
\end{definition}

\begin{theorem}[Largest certified bisimulation]\label{thm:largest-certified-bisimulation}
Under the closure condition of Proposition~\ref{prop:certified-subcoalgebra},
the relation \(\sim_{\cert}\) is itself a certified bisimulation on
\(X_{\cert}\) and is the largest such relation. Moreover:
\begin{enumerate}
  \item If \(x\sim_{\cert}y\), then \(\sem{x}_{\Gamma}=\sem{y}_{\Gamma}\).
  \item \(\sim_{\cert}\) is an equivalence relation on \(X_{\cert}\).
  \item Every raw bisimulation contained in
  \(X_{\cert}\times X_{\cert}\) is a certified bisimulation.
\end{enumerate}
\end{theorem}

\begin{proof}
Let \(\mathcal{B}\) be the family of all certified bisimulations and let
\(\sim_{\cert}=\bigcup_{R\in\mathcal{B}}R\). If
\((x,y)\in\sim_{\cert}\), then \((x,y)\in R\) for some
\(R\in\mathcal{B}\). Relation \(R\) supplies the matching successor in
both directions, and its successor pair belongs to
\(R\subseteq\sim_{\cert}\). Hence the union is a certified bisimulation;
maximality follows from its definition.

For~(1), suppose \((x,y)\in R\) for a certified bisimulation \(R\).
Because \(x,y\in X_{\cert}\), for every observation their unique raw
transitions are certified; closure keeps their successors in
\(X_{\cert}\). We prove by induction on the length of
\(u\in\Obs_{\Gamma}^{+}\) that
\(\sem{x}_{\Gamma}(u)=\sem{y}_{\Gamma}(u)\). For \(u=o\), write
\(x\overset{e,o}{\Longrightarrow}_{\cert}x'\). The first matching
condition produces \(y\overset{e,o}{\Longrightarrow}_{\cert}y'\), so the
characteristic equations give equal outputs. For \(u=ow\), the matched
successors satisfy \((x',y')\in R\); the induction hypothesis for \(w\)
and the characteristic equations give the result.

For~(2), the identity relation on \(X_{\cert}\) is a certified
bisimulation. If \(R\) is a certified bisimulation, then \(R^{-1}\) is one
because the definition contains both matching directions. If \(R\) and
\(S\) are certified bisimulations, their relational composite \(R\circ S\)
is one: match a step through \(R\) and then through \(S\), in either
direction. Since \(\sim_{\cert}\) contains the identity, inverses, and
composites of its witnesses, it is reflexive, symmetric, and transitive.

For~(3), let \(R\) be a raw bisimulation contained in
\(X_{\cert}\times X_{\cert}\). For \((x,y)\in R\) and observation \(o\),
total determinism gives unique transitions \(C(x)(o)=(e,x')\) and
\(C(y)(o)=(e',y')\). Raw bisimulation gives \(e=e'\) and
\((x',y')\in R\). Both transitions are certified because \(x,y\in
X_{\cert}\). The same matched pair therefore witnesses both directions in
Definition~\ref{def:certified-bisimulation}, so \(R\) is a certified
bisimulation.
\end{proof}

This gives a coinductive sufficient criterion for comparing certifiable prompt
decompositions, tool routings, or handoff policies by observable behaviour. The
claim is intentionally restricted to \(X_{\cert}\); the partial certified LTS
outside that subcoalgebra supports safety monitoring, but not the raw-behaviour
equality conclusion above.

\section{A SELL Layer for Agentic Resources}\label{sec:sell-layer}

Coalgebra describes what transitions are possible. It does not, by itself, say whether a transition is permitted, affordable, auditable, or resource-safe. For that we add a proof-theoretic layer.

\subsection{Subexponential signature}

Recall from Definition~\ref{def:subexp-signature} that a SELL signature is a triple
\[
  \Sigma=(I,\preceq,U),
\]
with subexponential labels \(I\), a preorder \(\preceq\) on \(I\), and an unbounded set \(U\subseteq I\) of labels admitting weakening and contraction; labels outside \(U\) are linear, so their assumptions cannot be freely copied or discarded. We now instantiate this signature for an agent swarm, using labels of the following form:
\[
  I_{\Gamma}=
  \{\mathsf{lin},\mathsf{prog},\mathsf{shared},\mathsf{mem},\mathsf{trace},\mathsf{obl}\}
  \cup
  \{\mathsf{agent}_i,\mathsf{policy}_i,\mathsf{tool}_i,\mathsf{budget}_i\mid i\in V\}.
\]
A minimal reusable set is
\[
  U_{\Gamma}=\{\mathsf{prog},\mathsf{shared},\mathsf{mem},\mathsf{trace}\}
  \cup\{\mathsf{agent}_i,\mathsf{policy}_i,\mathsf{tool}_i\mid i\in V\},
\]
while \(\mathsf{budget}_i\), \(\mathsf{lin}\), and \(\mathsf{obl}\) remain non-structural. Thus tool permissions and recorded traces may be reused as evidence, whereas budget tokens and trace obligations are linear.

\begin{table}[htbp]
\centering
\begin{tabular}{>{\raggedright\arraybackslash}p{0.24\linewidth}>{\raggedright\arraybackslash}p{0.54\linewidth}>{\raggedright\arraybackslash}p{0.14\linewidth}}
\toprule
Label & Intended meaning & Structural status \\
\midrule
\(\mathsf{mem}\) & Persistent memory facts or summaries & reusable \\
\(\mathsf{prog}\) & Reusable certificate-program clauses & reusable \\
\(\mathsf{shared}\) & Shared blackboard or public context & reusable \\
\(\mathsf{trace}\) & Emitted audit facts, usable as evidence & reusable \\
\(\mathsf{tool}_i\) & Reusable permission for tools available to \(i\) & reusable \\
\(\mathsf{policy}_i\) & Reusable local policy facts of agent \(i\) & reusable \\
\(\mathsf{agent}_i\) & Reusable local context of agent \(i\) & reusable \\
\(\mathsf{budget}_i\) & Consumable budget of calls, tokens, time, or money & linear \\
\(\mathsf{obl}\) & Obligation to emit an audit item & linear \\
\(\mathsf{lin}\) & Ordinary linear resources & linear \\
\bottomrule
\end{tabular}
\caption{A candidate subexponential discipline for tool-augmented agent swarms.}
\label{tab:subexponentials}
\end{table}

For the deterministic core, the preorder \(\preceq_{\Gamma}\) is the smallest reflexive and transitive relation satisfying, for every \(i\in V\),
\[
  \ell\preceq_{\Gamma}\mathsf{agent}_i
  \quad\text{for each }\ell\in
  \{\mathsf{prog},\mathsf{shared},\mathsf{mem},\mathsf{trace},\mathsf{lin},
  \mathsf{obl},\mathsf{policy}_i,\mathsf{tool}_i,\mathsf{budget}_i\},
\]
and, for every directed edge \((i,j)\in E_{\Gamma}\),
\[
  \mathsf{agent}_i\preceq_{\Gamma}\mathsf{agent}_j.
\]
We read \(\ell\preceq_{\Gamma}m\) as saying that resources stored at label \(\ell\) may be inspected by a rule focused at label \(m\), subject to the structural status of \(\ell\). Thus a rule for agent \(j\) may use shared facts, memory facts, trace evidence, its own tool permissions, and its own budget tokens, and additionally the local context of any agent \(i\) from which there is a directed path to \(j\) in \(E_{\Gamma}\). Agents not reachable into \(j\) have their private context inaccessible at the level of the SELL signature, independently of any runtime side condition. Proposition~\ref{prop:graph-dependent-preorder} below makes this precise.

\begin{proposition}[Graph-dependent structure of \(\Sigma_{\Gamma}\)]\label{prop:graph-dependent-preorder}
Let \(E_{\Gamma}^{+}\) denote the positive transitive closure of \(E_{\Gamma}\): the set of pairs \((i,j)\) with \(i\ne j\) for which there exists a directed path of at least one edge from \(i\) to \(j\) in \(E_{\Gamma}\). Then:
\begin{enumerate}
  \item For distinct \(i,j\in V\), \(\mathsf{agent}_i\preceq_{\Gamma}\mathsf{agent}_j\) if and only if \((i,j)\in E_{\Gamma}^{+}\).
  \item Two interaction graphs \(\Gamma=(V,E_{\Gamma})\) and \(\Gamma'=(V,E')\) with the same agent set satisfy \(\preceq_{\Gamma}=\preceq_{\Gamma'}\) if and only if \(E_{\Gamma}^{+}=E_{\Gamma'}^{+}\): the SELL preorder is uniquely determined by the reachability structure of \(\Gamma\), not by its individual edges.
  \item If \(E_{\Gamma'}^{+}\subsetneq E_{\Gamma}^{+}\), then \(\preceq_{\Gamma'}\subsetneq\preceq_{\Gamma}\): strictly fewer reachable agent pairs correspond to strictly reduced cross-agent context accessibility in the SELL signature.
  \item Under the focus convention of the certificate fragment, a rule focused
  on \(\mathsf{agent}_j\) may inspect
  \(\mathsf{agent}_i\)'s local zone \((i\ne j)\) if and only if
  \((i,j)\in E_{\Gamma}^{+}\).
  \item Let \(i\equiv_{\Gamma}j\) mean \(i=j\) or mutual positive reachability. The quotient of
  the agent-label preorder by
  \(\mathsf{agent}_i\equiv\mathsf{agent}_j\) iff both
  \(\mathsf{agent}_i\preceq_{\Gamma}\mathsf{agent}_j\) and
  \(\mathsf{agent}_j\preceq_{\Gamma}\mathsf{agent}_i\) is isomorphic to
  the reachability poset of the strongly connected components of \(\Gamma\).
\end{enumerate}
\end{proposition}

\begin{proof}
For~(1), the agent--agent generating pairs of \(\preceq_{\Gamma}\) are exactly \(\{(\mathsf{agent}_i,\mathsf{agent}_j)\mid(i,j)\in E_{\Gamma}\}\). Since \(\preceq_{\Gamma}\) is the smallest reflexive and transitive closure of these pairs together with the base--agent pairs, transitivity adds \((\mathsf{agent}_i,\mathsf{agent}_k)\) whenever there is a sequence of edges \((i,j_1),(j_1,j_2),\ldots,(j_m,k)\) in \(E_{\Gamma}\). By the definition of \(E_{\Gamma}^{+}\), this holds if and only if \((i,k)\in E_{\Gamma}^{+}\). The base labels satisfy \(\ell\preceq_{\Gamma}\mathsf{agent}_i\) for all \(i\in V\) independently of the graph, and no agent label precedes a base label, so the agent--agent pairs are exactly as stated.

For~(2), by~(1) the agent--agent part of \(\preceq_{\Gamma}\) is determined entirely by \(E_{\Gamma}^{+}\), and the base--agent part is graph-independent. Hence \(\preceq_{\Gamma}=\preceq_{\Gamma'}\) iff the agent--agent parts agree, which by~(1) holds iff \(E_{\Gamma}^{+}=E_{\Gamma'}^{+}\).

For~(3), if \(E_{\Gamma'}^{+}\subsetneq E_{\Gamma}^{+}\), there exists \((i,j)\in E_{\Gamma}^{+}\setminus E_{\Gamma'}^{+}\). By~(1), \(\mathsf{agent}_i\preceq_{\Gamma}\mathsf{agent}_j\) but \(\mathsf{agent}_i\not\preceq_{\Gamma'}\mathsf{agent}_j\), giving \(\preceq_{\Gamma'}\subsetneq\preceq_{\Gamma}\).

For~(4), in the SELL proof system with signature \(\Sigma_{\Gamma}\), using a formula at label \(\ell\) in a rule focused at label \(m\) requires \(\ell\preceq_{\Gamma}m\). Setting \(\ell=\mathsf{agent}_i\) and \(m=\mathsf{agent}_j\) with \(i\ne j\), the condition is exactly \(\mathsf{agent}_i\preceq_{\Gamma}\mathsf{agent}_j\), which by~(1) holds iff \((i,j)\in E_{\Gamma}^{+}\).

For~(5), two agent labels become equal in the quotient preorder exactly when
their vertices are mutually reachable, which is exactly membership in the same
strongly connected component. Ordering the resulting equivalence classes by
the quotient preorder gives \([i]\le [j]\) exactly when the condensation graph
has a path from the component of \(i\) to that of \(j\). This is the
reachability order of the condensation DAG, giving the claimed isomorphism.
\end{proof}

\begin{remark}
Proposition~\ref{prop:graph-dependent-preorder} establishes a precise, but
structural, coupling: direct edges constrain communication events, while graph
reachability constrains cross-agent context inspection in the certificate
fragment. The quotient result shows exactly how much topology the signature
retains: it forgets individual edges inside a reachability-equivalent graph but
retains the condensation order. This does not make the transition coalgebra a
SELL proof system; it makes the same topology parameter govern two distinct
layers whose consistency can be checked.
\end{remark}

\subsection{From SELL Sequents to Executable Certificates}\label{subsec:sell-sequents-certificates}

We use SELL as the proof-theoretic discipline for a transition certificate over
resource ledgers. To avoid conflating this use with theorem proving in full
SELL, we distinguish two objects. The \emph{certificate calculus} consists of
six side-conditioned operational schemas. Its \emph{SELL program
interpretation} places a reusable clause for each resource transformation in
the \(\mathsf{prog}\) zone. A schema instance is admissible only after its
non-logical side conditions---such as graph membership, arithmetic comparison,
payload policy, and freshness---have been checked. Conditional on that
evidence, focused use of the corresponding program clause consumes and produces
the displayed linear and reusable facts. We therefore claim an exact
checker--calculus correspondence and a resource-level SELL interpretation; we
do not claim that the displayed fractions are primitive rules of pure SELL or
that the checker decides unrestricted SELL provability.

A judgement has the form
\[
  \Delta;\Theta\vdash_{\Sigma_{\Gamma}} e:(y,R)\leadsto(y',R'),
\]
where \(\Delta\) is a multiset of linear tokens and \(\Theta\) is a labelled subexponential context. In the present calculus,
\[
\begin{aligned}
  \Delta={}&
  \Delta_{\mathsf{lin}}\uplus
  \Delta_{\mathsf{obl}}\uplus
  \biguplus_{i\in V}\Delta_{\mathsf{budget}_i},
  \qquad
  \Delta_{\mathsf{budget}_i}\ni \mathsf{budget}(i,n),\; n\in\Nat,\\
  \Theta={}&
  \Theta_{\mathsf{prog}}\uplus
  \Theta_{\mathsf{shared}}\uplus
  \Theta_{\mathsf{mem}}\uplus
  \Theta_{\mathsf{trace}}\uplus
  \biguplus_{i\in V}
  (\Theta_{\mathsf{agent}_i}\uplus\Theta_{\mathsf{policy}_i}\uplus
   \Theta_{\mathsf{tool}_i}).
\end{aligned}
\]
The reusable components of \(\Theta\) contain program clauses and formulas such as \(!^{\mathsf{tool}_i}\mathsf{mayUse}(i,t)\), \(!^{\mathsf{shared}}\mathsf{source}(k,v,p)\), \(!^{\mathsf{trace}}\mathsf{emitted}(q)\), and \(!^{\mathsf{mem}}\mathsf{remembered}(k)\). Linear components of \(\Delta\) contain consumable budget tokens and audit obligations. The bridge to the executable checker is obtained by reading each successful side-condition check and rule instance as a certificate object whose conclusion fixes the next ledger \(R'\). Failed premises become checker rejections.

The ledger-to-context interpretation used in the fragment is:
\[
\begin{aligned}
  \mathcal{I}(R,M)=&
  \{\mathsf{budget}(i,\beta(i))@\mathsf{budget}_i\mid i\in V\}\uplus
  \{!^{\mathsf{tool}_i}\mathsf{mayUse}(i,t)\mid (i,t)\in\pi\}\uplus{}\\
  &\{!^{\mathsf{trace}}\mathsf{emitted}(q)\mid q\in\tau\}\uplus
  \{!^{\mathsf{shared}}\mathsf{source}(k,v,p)\mid M(k)=(v,p)\}\uplus{}\\
  &\{!^{\mathsf{shared}}\mathsf{certified}(a,k)\mid (a,k)\in\kappa\}.
\end{aligned}
\]
The interpretation makes explicit why the calculus is resource-sensitive: permissions, source facts, trace facts, and certified-claim facts may be reused, but the budget token for \(i\) is a linear representative of the current numeric budget and is replaced rather than duplicated.

For a tool call by agent \(i\) using tool \(t\), let \(q\) be a fresh trace identifier, let \(r=\cost(t)\), and let \(W\) be a finite set of shared writes \((k,v,p)\) produced by the call. Every \(k\) in \(W\) must be fresh and every \(p\) must contain \((i,t,q)\), written \(\mathsf{prov}(p;i,t,q)\). Write \(M[W]\) for their simultaneous extension. The rule is sound only when \((i,t)\in\pi\), \(\beta(i)\ge r\), the write conditions hold, and
\[
  R'=(\beta[i\mapsto\beta(i)-r],\pi,\tau\cup\{q\},\kappa).
\]
At the certificate-calculus level this is represented by the schema
\[
\frac{
  \begin{gathered}
  !^{\mathsf{tool}_i}\mathsf{mayUse}(i,t)
  \quad \mathsf{budget}(i,\beta(i))@\mathsf{budget}_i
  \quad q\notin\tau \quad \beta(i)\ge \cost(t)\\
  \forall(k,v,p)\in W.\;M(k)\text{ undefined and }\mathsf{prov}(p;i,t,q)
  \end{gathered}
}{
  \begin{gathered}
  \Delta;\Theta\vdash_{\Sigma_{\Gamma}}\\
  \mathsf{call}(i,t,W,q):(y,R)\leadsto(y[M\mapsto M[W]],R')
  \end{gathered}
}.
\]
The budget premise is linear: the certificate consumes the old budget token and produces the reduced one. The write premises give append-only provenance, and trace freshness prevents one identifier from witnessing two obligations. We abbreviate \(\mathsf{call}(i,t,\emptyset,q)\) as \(\mathsf{call}(i,t,q)\); the JSON field \texttt{writes} supplies \(W\) in the executable schema.

A resource-level SELL reading of the same schema is the following transformer:
\[
\begin{aligned}
  &\mathsf{budget}(i,n)@\mathsf{budget}_i,
  !^{\mathsf{tool}_i}\mathsf{mayUse}(i,t),
  \mathsf{obl}(q)@\mathsf{obl}\ \vdash\\
  &\quad \mathsf{budget}(i,n-\cost(t))@\mathsf{budget}_i
  \otimes !^{\mathsf{trace}}\mathsf{emitted}(q)\\
  &\quad \otimes
  \bigotimes_{(k,v,p)\in W}!^{\mathsf{shared}}\mathsf{source}(k,v,p),
\end{aligned}
\]
with the side condition \(n\ge \cost(t)\) and \(q\notin\tau\). This transformer is the body of a reusable program clause at \(\mathsf{prog}\); it is not derivable from the resource facts alone. This is the precise sense in which the rule is SELL-indexed: the program and permission are reusable, the trace fact becomes reusable evidence, and the budget and trace obligation are linear.

If the tool call writes a shared-memory key \(k\) with value \(v\) and provenance \(p=(i,t,q)\), the memory-write schema requires freshness of the key and the corresponding audit item:
\[
\frac{
  M(k)\text{ undefined}
  \qquad
  q\notin\tau
  \qquad
  p=(i,t,q)
}{
  \Delta;\Theta\vdash_{\Sigma_{\Gamma}}
  \mathsf{writeShared}(i,k,v,p,q):(y,R)\leadsto(y[M\mapsto M[k\mapsto(v,p)]],R')
}.
\]
Here \(R'=(\beta,\pi,\tau\cup\{q\},\kappa)\). This rule is the proof-theoretic counterpart of the append-only order \(M\sqsubseteq M'\): once a key is available as source evidence, later transitions may cite it but may not silently overwrite it.

For a handoff from \(i\) to \(j\), soundness requires \((i,j)\in E_{\Gamma}\), a payload satisfying the handoff policy, a fresh trace item \(q\), and
\[
  R'=(\beta,\pi,\tau\cup\{q\},\kappa).
\]
The corresponding schema is
\[
\frac{
  (i,j)\in E_{\Gamma}
  \qquad
  \mathsf{policy}_{ij}(p)
  \qquad
  q\notin\tau
}{
  \Delta;\Theta\vdash_{\Sigma_{\Gamma}}
  \mathsf{handoff}_{ij}(p,q):(y,R)\leadsto(y',R')
}.
\]
Messages use the same graph-and-audit discipline, without transferring control:
\[
\frac{
  (i,j)\in E_{\Gamma}
  \qquad
  \mathsf{msgPolicy}_{ij}(m)
  \qquad
  q\notin\tau
}{
  \Delta;\Theta\vdash_{\Sigma_{\Gamma}}
  \mathsf{message}_{ij}(m,q):(y,R)\leadsto(y',R')
}.
\]
Here again \(R'=(\beta,\pi,\tau\cup\{q\},\kappa)\). For certification, a verifier may add a claim \(a\) supported by key \(k\) only when the key is already present in shared memory and the certifying agent is authorised:
\[
\frac{
  i\in\mathsf{Certifiers}
  \qquad
  M(k)=(v,p)
  \qquad
  q\notin\tau
}{
  \Delta;\Theta\vdash_{\Sigma_{\Gamma}}
  \mathsf{certify}(i,a,k,q):(y,R)\leadsto(y,R')
},
\]
where \(R'=(\beta,\pi,\tau\cup\{q\},\kappa\cup\{(a,k)\})\). Finally, an answer rule can cite only already certified claims:
\[
\frac{
  \forall a\in A\;\exists k_a.\,(a,k_a)\in\kappa
  \qquad
  q\notin\tau
}{
  \Delta;\Theta\vdash_{\Sigma_{\Gamma}}
  \mathsf{answer}(i,A,q):(y,R)\leadsto(y,R')
},
\]
where \(R'=(\beta,\pi,\tau\cup\{q\},\kappa)\). These six schemas are the minimal calculus used in the preservation theorem and in the prototype: tool calls spend budget, memory writes create provenance-carrying source keys, handoffs and messages enforce the graph policy, certification binds claims to those keys, and answers use only certified evidence.

The six schemas have a uniform SELL-indexed resource-transformer reading. Each displayed implication is stored as a reusable clause under \(\mathsf{prog}\). Formulas of the form \(!^{\ell}A\) are reusable when \(\ell\in U_{\Gamma}\), whereas budget tokens and audit obligations remain linear. The side conditions are exactly those stated in the operational schemas: known endpoints, graph edges, payload policies, source-key freshness, trace freshness, certifier authority, source availability, and budget sufficiency.
\[
\begin{array}{@{}r@{\quad}l@{}}
\mathsf{call}:&
\begin{gathered}[t]
  \mathsf{budget}(i,n)@\mathsf{budget}_i,
  !^{\mathsf{tool}_i}\mathsf{mayUse}(i,t),
  \mathsf{obl}(q)@\mathsf{obl}\vdash{}\\
  \mathsf{budget}(i,n-\cost(t))@\mathsf{budget}_i
  \otimes \displaystyle\bigotimes_{(k,v,p)\in W}
    !^{\mathsf{shared}}\mathsf{source}(k,v,p)
  \otimes !^{\mathsf{trace}}\mathsf{emitted}(q)
\end{gathered}\\[2mm]
\mathsf{writeShared}:&
\begin{gathered}[t]
  \mathsf{freshKey}(k),
  \mathsf{prov}(p,i,t,q),
  \mathsf{obl}(q)@\mathsf{obl}\vdash{}\\
  !^{\mathsf{shared}}\mathsf{source}(k,v,p)
  \otimes !^{\mathsf{trace}}\mathsf{emitted}(q)
\end{gathered}\\[2mm]
\mathsf{handoff}:&
\begin{gathered}[t]
  !^{\mathsf{shared}}\mathsf{edge}_{\Gamma}(i,j),
  !^{\mathsf{policy}_i}\mathsf{handoffPolicy}_{ij}(p),
  \mathsf{obl}(q)@\mathsf{obl}\vdash{}\\
  !^{\mathsf{trace}}\mathsf{handoff}(i,j,q)
  \otimes !^{\mathsf{trace}}\mathsf{emitted}(q)
\end{gathered}\\[2mm]
\mathsf{message}:&
\begin{gathered}[t]
  !^{\mathsf{shared}}\mathsf{edge}_{\Gamma}(i,j),
  !^{\mathsf{policy}_i}\mathsf{msgPolicy}_{ij}(m),
  \mathsf{obl}(q)@\mathsf{obl}\vdash{}\\
  !^{\mathsf{trace}}\mathsf{message}(i,j,q)
  \otimes !^{\mathsf{trace}}\mathsf{emitted}(q)
\end{gathered}\\[2mm]
\mathsf{certify}:&
\begin{gathered}[t]
  !^{\mathsf{policy}_i}\mathsf{mayCertify}(i),
  !^{\mathsf{shared}}\mathsf{source}(k,v,p),
  \mathsf{obl}(q)@\mathsf{obl}\vdash{}\\
  !^{\mathsf{shared}}\mathsf{certified}(a,k)
  \otimes !^{\mathsf{trace}}\mathsf{emitted}(q)
\end{gathered}\\[2mm]
\mathsf{answer}:&
\begin{gathered}[t]
  \bigotimes_{a\in A}!^{\mathsf{shared}}\mathsf{certified}(a,k_a),
  \mathsf{obl}(q)@\mathsf{obl}\vdash{}\\
  !^{\mathsf{trace}}\mathsf{answered}(A,q)
  \otimes !^{\mathsf{trace}}\mathsf{emitted}(q)
\end{gathered}.
\end{array}
\]
Formally, let \(\mathcal{P}_{\mathcal{S}}\) contain the universal closures of
these six implications, stored as
\(!^{\mathsf{prog}}(L\multimap R)\), for the finite specification
\(\mathcal{S}\). An admissible ground instance may be focused only when its
external side-condition evidence is present. One focused use consumes the
linear formulas in \(L\), retains reusable premises, and produces the formulas
in \(R\). The cut between the certification conclusion and an answer premise is
permitted because certified-claim facts live in the reusable shared zone. By
contrast, each audit obligation remains linear, so the same obligation cannot
justify two auditable steps; the separate freshness side condition ensures that
their identifiers are distinct in the ledger.

\subsection{Status of the SELL fragment}\label{subsec:sell-status}

The relationship between the present calculus and full SELL proof theory is
governed by three explicit boundaries.

\begin{enumerate}
\item \textbf{No complete proof-search.}
  The six schemas do not implement a complete proof-search procedure for SELL.
  There is no focusing strategy, no backtracking, and no general sequent-theorem prover for arbitrary SELL formulae.
  Each schema is a fixed rule template whose premises are directly checkable against the current runtime ledger.

\item \textbf{Side-conditioned SELL program fragment.}
  Each admissible rule instance is one focused use of a clause in
  \(\mathcal{P}_{\mathcal{S}}\). Reusable assumptions (program clauses, tool
  permissions, source facts, certified claims, and trace evidence) occur under
  labels in \(U_{\Gamma}\), while budget tokens and audit obligations occur in
  the linear context \(\Delta\). Arithmetic, freshness, and application policy
  remain explicit side conditions rather than being claimed as theorems of pure
  SELL.
  The fragment is \emph{finite} in the sense that the set of possible rule instances for a given specification
  \(\mathcal{S}\) is determined entirely by the finite sets of agents, tools, graph edges, certifiers, and
  budget constants in \(\mathcal{S}\).

\item \textbf{Decidability.}
  Because each premise reduces to a membership or freshness test over sets derived from the finite
  specification \(\mathcal{S}\), the fragment is decidable.
  The decision procedure is the abstract checker \(\TraceCheck_{\mathcal{S}}\) defined in
  Section~\ref{sec:prototype} (Definition~\ref{def:tracecheck} and Algorithm~\ref{alg:tracecheck}):
  it processes a finite event sequence from left to right in a single pass, performing only
  finite membership and freshness tests at each step.
\end{enumerate}

These boundaries support the correspondence result
(Proposition~\ref{prop:checker-adequacy}): accepted traces are exactly those
generated by admissible instances of the six schemas. The proposition is an
executable-specification result, not a new completeness theorem for SELL.
Section~\ref{subsec:focusing-adequacy} sharpens the positive side of these
boundaries: within the fixed program \(\mathcal{P}_{\mathcal{S}}\), each
admissible instance is exactly one focused phase
(Theorem~\ref{thm:focusing-adequacy}).

\begin{definition}[Local soundness of certificates]\label{def:local-soundness}
A family of SELL transition rules is locally sound for \(\Gamma\) if each certified step
\[
  e:(y,R)\leadsto(y',R')
\]
with shared-memory components \(M,M'\), \(R=(\beta,\pi,\tau,\kappa)\), and \(R'=(\beta',\pi',\tau',\kappa')\) satisfies the following conditions:
\begin{enumerate}
  \item graph condition: if \(e\) is a message or handoff from \(i\) to \(j\), then \((i,j)\in E_{\Gamma}\);
  \item payload-policy condition: every message or handoff payload satisfies the configured policy for its directed edge;
  \item permission condition: if \(e\) is a tool call by \(i\) to tool \(t\), then \((i,t)\in\pi\);
  \item permission-ledger condition: \(\pi'=\pi\) in the present calculus, which has no implicit permission grant or revocation;
  \item shared-memory condition: \(M\sqsubseteq M'\), and every newly allocated key in \(\operatorname{dom}(M')\setminus\operatorname{dom}(M)\) carries provenance containing the event and its trace item;
  \item budget condition: for each agent \(i\),
  \[
    c_i(e)\le \beta(i)+\rho_i(e)
    \qquad\text{and}\qquad
    \beta'(i)=\beta(i)+\rho_i(e)-c_i(e),
  \]
  where \(c_i(e),\rho_i(e)\in\Nat\) are the certified cost and an explicitly certified replenishment; the inequality makes the subtraction an ordinary non-negative integer subtraction rather than truncated subtraction in \(\Nat\);
  \item trace condition: \(\tau\subseteq\tau'\), and if \(e\) is auditable, then there exists a trace item \(q_e\in\tau'\setminus\tau\);
  \item certification monotonicity: \(\kappa\subseteq\kappa'\);
  \item certification condition: every pair \((a,k)\in\kappa'\setminus\kappa\) is supported by a source key \(k\in\operatorname{dom}(M)\) already present before the certification step;
  \item answer condition: if \(e\) is a final answer citing claims \(K_e\), then for every \(a\in K_e\) there is some key \(k\) such that \((a,k)\in\kappa\).
\end{enumerate}
\end{definition}


\begin{table}[htbp]
\centering
\small
\begin{tabular}{>{\raggedright\arraybackslash}p{0.22\linewidth}>{\raggedright\arraybackslash}p{0.32\linewidth}>{\raggedright\arraybackslash}p{0.34\linewidth}}
\toprule
Operational rule & SELL-labelled resources & Local obligation discharged \\
\midrule
\(\mathsf{call}(i,t,W,q)\) & \(!^{\mathsf{tool}_i}\mathsf{mayUse}(i,t)\), \(\mathsf{budget}(i,n)@\mathsf{budget}_i\), fresh write and trace obligations & permission, budget decrement, append-only provenance, fresh audit item \\
\(\mathsf{writeShared}(i,k,v,p,q)\) & shared-memory source key, provenance term \(p\), fresh trace obligation & append-only memory extension, provenance recording, fresh audit item \\
\(\mathsf{handoff}_{ij}(p,q)\) & agent-local payload policy, graph fact, fresh trace obligation & graph preservation, payload admissibility, fresh audit item \\
\(\mathsf{message}_{ij}(m,q)\) & message policy, graph fact, fresh trace obligation & graph preservation, message admissibility, fresh audit item \\
\(\mathsf{certify}(i,a,k,q)\) & certifier authority, \(!^{\mathsf{shared}}\mathsf{source}(k,v,p)\), fresh trace obligation & claim-source binding, certification monotonicity, fresh audit item \\
\(\mathsf{answer}(i,A,q)\) & certified claim-source pairs \((a,k_a)\in\kappa\), fresh trace obligation & answer safety and traceability \\
\bottomrule
\end{tabular}
\caption{Executable certificate rules and the local soundness obligations they discharge.}
\label{tab:certificate-obligations}
\end{table}

\begin{lemma}[Local soundness of the certification calculus]\label{lem:calculus-local-soundness}
Every instance of the six certificate schemas above that satisfies its stated premises satisfies the corresponding clauses of Definition~\ref{def:local-soundness}.
\end{lemma}

\begin{proof}
The proof is by case analysis on the last rule used to derive the certificate. For \(\mathsf{call}(i,t,W,q)\), the permission and budget premises give \((i,t)\in\pi\) and \(\beta(i)\ge\cost(t)\), and the conclusion subtracts exactly that cost while leaving other budgets and the permission ledger unchanged. Fresh, provenance-carrying writes give \(M\sqsubseteq M[W]\). The side condition \(q\notin\tau\) and update \(\tau' = \tau\cup\{q\}\) give trace freshness; certified claims are unchanged. Hence all corresponding local-soundness clauses hold.

For \(\mathsf{writeShared}(i,k,v,p,q)\), freshness of \(k\) and the update \(M'=M[k\mapsto(v,p)]\) give \(M\sqsubseteq M'\). The premise fixing \(p=(i,t,q)\), or more generally a provenance term containing the event and trace item, gives the shared-memory provenance clause. Again \(q\notin\tau\) and \(\tau' = \tau\cup\{q\}\) give trace freshness, while budgets and certified claims are unchanged.

For \(\mathsf{handoff}_{ij}(p,q)\), the premise \((i,j)\in E_{\Gamma}\) is exactly the graph condition. The payload-policy premise states that the handoff payload is admissible. The trace update gives a fresh audit item, and no budget, memory, or certification component is changed.

For \(\mathsf{message}_{ij}(m,q)\), the same graph premise gives the graph condition for messages, the message-policy premise states admissibility of the message payload, and the trace update gives a fresh audit item. Budgets, memory, and certified claims are unchanged.

For \(\mathsf{certify}(i,a,k,q)\), the certifier-authority premise authorises the rule instance, and the premise \(M(k)=(v,p)\) ensures that the source key is already present before the claim is added. The conclusion updates \(\kappa\) to \(\kappa\cup\{(a,k)\}\), so certification is monotone and every newly added claim has a shared-memory source. Trace freshness follows from the same side condition and update as above.

For \(\mathsf{answer}(i,A,q)\), the premise \(\forall a\in A\,\exists k_a.(a,k_a)\in\kappa\) is exactly the answer-safety condition. The conclusion leaves budgets, memory, and certifications unchanged and adds the fresh audit item \(q\). These cases exhaust the calculus, proving the lemma.
\end{proof}

The purpose of Lemma~\ref{lem:calculus-local-soundness} is not to establish a new metatheorem of SELL. Its role is \emph{architectural}: it is the precise bridge between the six labelled rule schemas of Section~\ref{subsec:sell-sequents-certificates} and the global resource-preservation statement of Theorem~\ref{thm:resource-preservation}. Without it, the preservation theorem would depend on informally assumed rule correctness. With it, the proof obligations are finitely enumerated and each case is independently verifiable. This structure mirrors compiler correctness proofs~\cite{leroy2009compcert}, where local instruction-preservation lemmas are stated case by case before they are composed into a global simulation argument. The explicit case decomposition also delineates the mechanisation target: encoding each case as a goal in Coq or Agda would replace the present paper proof with a mechanically checked one, upgrading the bridge from informal to formal without changing the statement of Theorem~\ref{thm:resource-preservation}.

\subsection{Focused Adequacy of the Certificate Fragment}\label{subsec:focusing-adequacy}

The three boundaries above say what the fragment is \emph{not}. This subsection
states, positively, the exact proof-theoretic status of the six schemas: modulo
external side-condition evidence, each admissible rule instance is one complete
focused phase over a bipolar clause of \(\mathcal{P}_{\mathcal{S}}\), so the
checker recognises derivability in focused SELL from \(\mathcal{P}_{\mathcal{S}}\)
at the level of derivations in the sense of Nigam and
Miller~\cite{nigam2009,nigammiller2010}. We first isolate the syntactic shape
that makes focusing collapse each rule to a single phase.

Recall from Section~\ref{subsec:sell-sequents-certificates} that every clause of
\(\mathcal{P}_{\mathcal{S}}\) is stored as \(!^{\mathsf{prog}}(L\multimap R)\),
where \(L\) and \(R\) are finite tensor products of labelled atoms: the linear
tokens \(\mathsf{budget}(i,n)@\mathsf{budget}_i\) and
\(\mathsf{obl}(q)@\mathsf{obl}\), together with reusable atoms \(!^{\ell}A\) for
\(\ell\in U_{\Gamma}\), namely \(!^{\mathsf{tool}_i}\mathsf{mayUse}\),
\(!^{\mathsf{shared}}\mathsf{source}\),
\(!^{\mathsf{shared}}\mathsf{certified}\),
\(!^{\mathsf{shared}}\mathsf{edge}_{\Gamma}\), \(!^{\mathsf{policy}_i}(\cdot)\),
and \(!^{\mathsf{trace}}(\cdot)\). We use a standard focused SELL proof
system~\cite{andreoli1992,nigam2009}: a derivation alternates negative
(asynchronous) and positive (synchronous) phases, and once a formula is selected
for focus, focusing persists on its positive subformulae until a negative
formula or atom is reached.

\begin{definition}[Bipolar certificate clause]\label{def:bipolar}
A clause \(!^{\mathsf{prog}}(L\multimap R)\) is \emph{bipolar} when \(L\) and
\(R\) are tensors of labelled atoms of the form above, with no additive
connective, no implication nested inside \(L\) or \(R\), and no quantifier
alternation beyond the outer universal closure. Focusing on a bipolar clause
therefore comprises exactly one positive phase that consumes the atoms of \(L\)
and releases the atoms of \(R\), with no intervening decision.
\end{definition}

\begin{proposition}[Bipolarity of the fragment]\label{prop:bipolarity}
Every clause of \(\mathcal{P}_{\mathcal{S}}\) for the six schemas
\(\mathsf{call}\), \(\mathsf{writeShared}\), \(\mathsf{handoff}\),
\(\mathsf{message}\), \(\mathsf{certify}\), and \(\mathsf{answer}\) is bipolar.
\end{proposition}

\begin{proof}
Inspecting the six transformers in
Section~\ref{subsec:sell-sequents-certificates}, each left-hand side \(L\) is a
tensor of budget, obligation, permission, edge, policy, source, and certified
atoms, and each right-hand side \(R\) is a tensor of budget, trace, source, and
certified atoms. No clause contains an additive connective, an implication
nested inside \(L\) or \(R\), or a quantifier alternation beyond the outer
universal closure. Hence each clause matches Definition~\ref{def:bipolar}.
\end{proof}

\begin{theorem}[Focused adequacy modulo side conditions]\label{thm:focusing-adequacy}
Fix \(\Sigma_{\Gamma}\), the program \(\mathcal{P}_{\mathcal{S}}\), and the
ledger interpretation \(\mathcal{I}(R,M)\) of
Section~\ref{subsec:sell-sequents-certificates}. For each of the six schemas and
each ground event \(e\) whose external side conditions---graph membership,
payload policy, arithmetic \(\beta(i)\ge\cost(t)\), and freshness
\(q\notin\tau\)---hold at \((R,M)\), the following are in one-to-one
correspondence:
\begin{enumerate}
  \item an admissible instance of the schema \(e:(y,R)\leadsto(y',R')\) with
  shared memory \(M\mapsto M'\);
  \item one focused phase that decides on the corresponding bipolar clause of
  \(\mathcal{P}_{\mathcal{S}}\) and rewrites the labelled context
  \(\mathcal{I}(R,M)\) to \(\mathcal{I}(R',M')\).
\end{enumerate}
Consequently, a finite event sequence \(e_0\cdots e_{n-1}\) with all side
conditions met is generated by admissible schema instances producing ledgers
\(R_0,\dots,R_n\) if and only if there is a focused SELL derivation from
\(\mathcal{P}_{\mathcal{S}}\) and \(\mathcal{I}(R_0,M_0)\) that reaches
\(\mathcal{I}(R_n,M_n)\) by \(n\) successive phases, one per event.
\end{theorem}

\begin{proof}
Fix a schema and a ground event \(e\) whose side conditions hold. By
Proposition~\ref{prop:bipolarity} the corresponding clause
\(!^{\mathsf{prog}}(L\multimap R)\) is bipolar, and \(\mathsf{prog}\in U_{\Gamma}\)
makes it available for focus without consuming it. Deciding on the clause opens a
single positive phase. The phase consumes exactly the linear atoms of \(L\)
present in \(\mathcal{I}(R,M)\)---the budget token \(\mathsf{budget}(i,n)\) and,
where present, the obligation \(\mathsf{obl}(q)\)---and matches the reusable
atoms of \(L\) against the structural zones, which is permitted because those
labels lie in \(U_{\Gamma}\) and satisfy the preorder condition of
Section~\ref{sec:sell-layer}. Matching the reusable premises succeeds precisely
when the corresponding membership side conditions hold, while the arithmetic and
freshness side conditions are the external guards required before the phase may
fire. The phase then releases the atoms of \(R\): the decremented budget token
\(\mathsf{budget}(i,n-\cost(t))\), the new reusable facts
\(!^{\mathsf{shared}}\mathsf{source}\), \(!^{\mathsf{shared}}\mathsf{certified}\),
or the appropriate \(!^{\mathsf{trace}}(\cdot)\) event fact, and the reusable
trace evidence \(!^{\mathsf{trace}}\mathsf{emitted}(q)\). By the ledger update
of the schema, the resulting context is exactly \(\mathcal{I}(R',M')\). This
establishes the correspondence \((1)\Leftrightarrow(2)\) for a single event: an
admissible instance determines a phase, and conversely a completed phase
deciding on the clause fixes the same \(R'\) and \(M'\).

For the sequence statement, compose the single-event correspondences. Reusable
outputs stay available across later phases because their labels admit
contraction, so a \(\mathsf{source}\) or \(\mathsf{certified}\) fact produced by
one phase may be matched by the \(\mathsf{certify}\) or \(\mathsf{answer}\) phase
of a later event. Linear budget and obligation tokens, by contrast, are consumed
by the phase that produces the next ledger, so no phase can reuse a spent token.
Induction on \(n\) yields the stated derivation-level equivalence between
admissible schema runs and focused phase sequences.
\end{proof}

\begin{corollary}[Cut admissibility for certificate composition]\label{cor:cut-admissibility}
Along any certified run, the cut that lets a later \(\mathsf{certify}\) or
\(\mathsf{answer}\) phase use a \(\mathsf{source}\) or \(\mathsf{certified}\)
fact produced by an earlier phase is admissible. No cut against a linear budget
or obligation token is admissible, so a spent budget or a discharged obligation
cannot justify a second step.
\end{corollary}

\begin{proof}
Facts \(!^{\mathsf{shared}}\mathsf{source}(k,v,p)\) and
\(!^{\mathsf{shared}}\mathsf{certified}(a,k)\) carry a label in \(U_{\Gamma}\),
which admits weakening and contraction; the standard SELL cut-elimination for
unbounded labels~\cite{nigam2009,despeyroux2017} makes cut against such formulae
admissible, so the cut formula persists for later phases. Budget and obligation
tokens carry labels outside \(U_{\Gamma}\); a cut duplicating them would violate
linearity, and the schemas never introduce such a cut.
\end{proof}

\begin{remark}
Theorem~\ref{thm:focusing-adequacy} is an adequacy statement at the level of
derivations, not a completeness theorem for arbitrary SELL formulae. Its scope is
the fixed program \(\mathcal{P}_{\mathcal{S}}\) and its bipolar clauses; the
arithmetic, graph-membership, payload-policy, and freshness predicates remain
external evidence rather than SELL theorems, exactly as recorded in the
boundaries above. What the theorem adds is the precise sense in which the
executable checker is a focused proof-search engine for this fragment: each
accepted event is one focused phase, and Corollary~\ref{cor:cut-admissibility}
explains why reusable certified evidence composes while linear budgets do not.
\end{remark}

\subsection{Resource preservation}

\begin{theorem}[Resource preservation for certified runs]\label{thm:resource-preservation}
Let
\[
  (y_0,R_0) \xrightarrow{e_0} (y_1,R_1)
  \xrightarrow{e_1}\cdots
  \xrightarrow{e_{n-1}} (y_n,R_n)
\]
be a finite execution of a swarm coalgebra such that every step is certified by a locally sound SELL rule. Write \(R_k=(\beta_k,\pi_k,\tau_k,\kappa_k)\), and let \(M_k\) be the shared-memory component of \(y_k\). Then for every agent \(i\in V\),
\[
  \sum_{k=0}^{n-1} c_i(e_k)
  \le
  \beta_0(i)+\sum_{k=0}^{n-1}\rho_i(e_k).
\]
Moreover, every message and handoff respects its graph edge and payload policy, shared memory grows by provenance-preserving extension, every auditable tool call, handoff, message, shared-memory update, certification, or answer in the run has a corresponding fresh trace item in \(\tau_n\), and every claim used by a final answer belongs to the certified-claim component available at the step where the answer is emitted and therefore to \(\kappa_n\) with a supporting shared-memory key.
\end{theorem}

\begin{proof}
The graph, payload-policy, and permission conclusions hold at each relevant
step directly by local soundness. It remains to show that the ledger properties
accumulate across the run.

Fix an agent \(i\in V\). By the budget condition in Definition~\ref{def:local-soundness}, every certified step satisfies
\[
  \beta_{k+1}(i)=\beta_k(i)+\rho_i(e_k)-c_i(e_k).
\]
We prove by induction on \(m\), for \(0\le m\le n\), that
\[
  \beta_m(i)=\beta_0(i)+\sum_{k=0}^{m-1}\rho_i(e_k)-\sum_{k=0}^{m-1}c_i(e_k).
\]
For \(m=0\), both sums are empty, so the right-hand side is \(\beta_0(i)\), which equals the left-hand side \(\beta_0(i)\). Assume the equality holds for \(m\). Then
\[
\begin{aligned}
  \beta_{m+1}(i)
  &=\beta_m(i)+\rho_i(e_m)-c_i(e_m)\\
  &=\beta_0(i)+\sum_{k=0}^{m-1}\rho_i(e_k)-\sum_{k=0}^{m-1}c_i(e_k)+\rho_i(e_m)-c_i(e_m)\\
  &=\beta_0(i)+\sum_{k=0}^{m}\rho_i(e_k)-\sum_{k=0}^{m}c_i(e_k).
\end{aligned}
\]
Thus the equality holds for all \(m\), in particular for \(m=n\). Since every ledger budget has codomain \(\Nat\) and the local sufficiency inequality makes each subtraction non-negative, \(\beta_n(i)\ge0\). Rearranging the equality for \(m=n\) yields
\[
  \sum_{k=0}^{n-1} c_i(e_k)
  \le
  \beta_0(i)+\sum_{k=0}^{n-1}\rho_i(e_k).
\]
This proves the budget preservation claim.

For shared memory, local soundness gives \(M_k\sqsubseteq M_{k+1}\) at every step. By induction on \(k\), \(M_0\sqsubseteq M_n\). If a key first appears at step \(k\), the shared-memory condition records provenance containing the event and its trace item; extension never changes that binding. Thus every source key available at the end of the run has stable provenance back to the step that introduced it.

For traces, let \(k\) be a step at which \(e_k\) is auditable. By local soundness, there exists \(q_k\in\tau_{k+1}\setminus\tau_k\). The trace condition also gives \(\tau_k\subseteq\tau_{k+1}\) for every step, so by induction on the suffix length,
\[
  \tau_{k+1}\subseteq\tau_n.
\]
Hence \(q_k\in\tau_n\). If \(k<\ell\) are two auditable steps and \(q_k=q_\ell\), then \(q_k\in\tau_{\ell}\) by monotonicity, contradicting \(q_\ell\in\tau_{\ell+1}\setminus\tau_{\ell}\). Therefore different auditable events receive distinct trace items.

For final answers, suppose \(e_k\) is an answer citing the set of claims \(K_{e_k}\). The answer condition gives, for every \(a\in K_{e_k}\), a key \(k_a\) such that \((a,k_a)\in\kappa_k\). Certification monotonicity gives \(\kappa_k\subseteq\kappa_n\) by induction over the remaining steps. Therefore \((a,k_a)\in\kappa_n\) for every cited claim. This proves the final-answer claim. The certification condition additionally ensures that any claim newly entering \(\kappa\) during the run is backed by a source in the corresponding shared-memory state. By the shared-memory argument, that source has stable provenance. Hence certified answers cannot rely on claims introduced without source support.
\end{proof}

\begin{corollary}[Prefix safety]\label{cor:prefix-safety}
If an infinite execution has the property that every finite prefix is certified by locally sound SELL rules, then every finite prefix satisfies the budget, graph, payload-policy, permission, provenance, certification, and audit conclusions of Theorem~\ref{thm:resource-preservation}.
\end{corollary}

\begin{proof}
Apply Theorem~\ref{thm:resource-preservation} to an arbitrary finite prefix. Since the prefix was arbitrary, the property holds for all finite prefixes.
\end{proof}

For a fixed raw coalgebra \(C\) and initial state \(x_0=(y_0,R_0)\), let \(\mathcal{L}_{\mathrm{raw}}(C,x_0)\) be the set of finite words
\[
  (o_0,e_0)\cdots(o_{n-1},e_{n-1})
\]
generated by raw transitions of \(C\). Let \(\mathcal{L}_{\cert}(C,\Sigma_{\Gamma},x_0)\) be the subset generated by the certified transition relation \(\Longrightarrow_{\cert}\) from the same initial state.

\begin{proposition}[Certified safety sublanguage]\label{prop:certified-safety-sublanguage}
For every \(C\), \(\Sigma_{\Gamma}\), and initial state \(x_0\),
\[
  \mathcal{L}_{\cert}(C,\Sigma_{\Gamma},x_0)
  \subseteq
  \mathcal{L}_{\mathrm{raw}}(C,x_0).
\]
Moreover, \(\mathcal{L}_{\cert}(C,\Sigma_{\Gamma},x_0)\) is closed under finite prefixes, certified paths compose whenever the terminal state and ledger of the first path are the initial state and ledger of the second, and every certified word satisfies the budget, provenance, certification, and audit conclusions of Theorem~\ref{thm:resource-preservation}.
\end{proposition}

\begin{proof}
The inclusion follows immediately from the definition of \(\Longrightarrow_{\cert}\): a certified step exists only when the corresponding raw coalgebra step exists and its event has an admissible certificate. Prefix closure follows because every prefix of a finite sequence of certified steps is again a finite sequence of certified steps beginning at \(x_0\). Composition is concatenation of certificate sequences: if the final state and ledger of the first sequence match the initial state and ledger of the second, the conclusions of the first provide exactly the premises for starting the second. Finally, each certified word is a finite certified run, so Theorem~\ref{thm:resource-preservation} applies to it. Thus certification defines a safety sublanguage of raw behaviour rather than an unrelated execution model.
\end{proof}

\begin{remark}
The theorem is conditional on explicit local soundness assumptions. This is the appropriate level of abstraction for the article: a full proof-search implementation of SELL is not required, but each executable transition rule must be shown to satisfy the stated resource, graph, permission, provenance, trace, and certification invariants. Corollary~\ref{cor:prefix-safety} and Proposition~\ref{prop:certified-safety-sublanguage} also clarify the role of finiteness: the proved theorem is finite-run, but the safety property it establishes is prefix-closed and therefore meaningful for long-running systems monitored one finite prefix at a time.

Proposition~\ref{prop:certified-safety-sublanguage} collects three useful closure properties. First, the certified sublanguage is \emph{prefix-closed}: any truncation of a certified word remains certified. Second, certified paths \emph{compose} when their boundary state and ledger agree. Third, rejected prefixes act as finite bad prefixes in the safety-monitoring sense of Alpern and Schneider~\cite{alpern1987safety}. These properties explain why the side-conditioned certificate calculus supports incremental monitoring; they are not presented as a new characterisation theorem for safety languages in general.
\end{remark}

\subsection{Separated Swarm Composition}

The resource-preservation theorem is proved for a fixed interaction graph. A
composition theorem requires more than disjoint agent names: audit identifiers,
memory keys, claims, and mutable state must also be protected from accidental
collision or cross-component reads. We therefore state the exact
non-interference condition under which local certificates compose.

\begin{definition}[Separated swarm composition]\label{def:swarm-composition}
Let \(\Gamma_k=(V_k,E_k)\), \(k=1,2\), with
\(V_1\cap V_2=\emptyset\). A \emph{separated composition} consists of:
\begin{enumerate}
  \item the graph
  \(\Gamma_1\oplus\Gamma_2=(V_1\uplus V_2,E_1\uplus E_2)\), with no
  cross-component edge;
  \item tagged, disjoint namespaces
  \(\Key=\Key_1\uplus\Key_2\),
  \(\Trace=\Trace_1\uplus\Trace_2\), and
  \(\Claim=\Claim_1\uplus\Claim_2\), where component \(k\) allocates and
  cites only its own tags;
  \item product memory and environment
  \(\Mem=\Mem_1\times\Mem_2\) and
  \(\Env=\Env_1\times\Env_2\), with a component step reading and writing
  only its own factor;
  \item a combined ledger
  \[
    R_1\oplus R_2=
    (\beta_1\sqcup\beta_2,\,\pi_1\uplus\pi_2,\,
     \iota_1\tau_1\cup\iota_2\tau_2,\,
     \iota_1\kappa_1\cup\iota_2\kappa_2),
  \]
  where the injections tag trace identifiers, claims, and source keys.
\end{enumerate}
For \(X=X_{\Gamma_1}\times X_{\Gamma_2}\), observations and events are
disjoint sums. The coalgebra \(C_1\otimes C_2\) acts by
\[
\begin{aligned}
  (C_1\otimes C_2)(x_1,x_2)(\iota_1 o)
    &= (\iota_1 e,(x_1',x_2))
       &&\text{when }C_1(x_1)(o)=(e,x_1'),\\
  (C_1\otimes C_2)(x_1,x_2)(\iota_2 o)
    &= (\iota_2 e,(x_1,x_2'))
       &&\text{when }C_2(x_2)(o)=(e,x_2').
\end{aligned}
\]
The ledger displayed above is the canonical identification of the two ledger
factors with the ledger component of this product state.
\end{definition}

\begin{theorem}[Compositionality under separation]\label{thm:swarm-compositionality}
Let each \((\Gamma_k,C_k,{\vdash}_{\Sigma_{\Gamma_k}})\) be locally sound,
and assume the separation conditions of Definition~\ref{def:swarm-composition}.
Then every interleaving of certified component runs is a certified run of
\(C_1\otimes C_2\), and it satisfies Theorem~\ref{thm:resource-preservation}
for the combined ledger. In particular:
\begin{enumerate}
  \item each tool call decreases only its component budget by its declared
  cost \(\cost(t)\);
  \item permissions are component-local, and each memory factor remains an
  append-only provenance map;
  \item tagged audit identifiers remain globally fresh, and tagged
  claim--source pairs remain supported in their component memory.
\end{enumerate}
\end{theorem}

\begin{proof}
Consider one step of an interleaving, from component \(k\). By the definition
of \(C_1\otimes C_2\), the other state factor is unchanged. The component's
certificate uses only its own agent labels and its tagged resource context.
Because \(E_1\uplus E_2\) has no cross-component path,
Proposition~\ref{prop:graph-dependent-preorder} introduces no cross-agent
accessibility between the components. The component rule therefore remains an
admissible instance of the combined certificate program.

Local soundness supplies exact budget accounting, permission checking,
append-only memory, claim support, and freshness within component \(k\). The
tag injections preserve freshness and disjointness globally: an identifier or
key allocated by one component cannot equal one allocated by the other. The
unchanged factor preserves the other component's invariants. Induction on the
length of the interleaving proves that every step is certified and that all
local-soundness clauses hold for the combined ledger. Applying
Theorem~\ref{thm:resource-preservation} gives the result.
\end{proof}

\begin{corollary}[Closure under separated composition]\label{cor:composition-closure}
Locally sound certified swarms are closed under the separated composition of
Definition~\ref{def:swarm-composition}.
\end{corollary}

\begin{remark}
Separation is a real hypothesis, not bookkeeping. If components share an
untagged memory pool, reuse audit identifiers, inspect one another's shared
facts, or add interface edges, their local certificates need not remain valid
without additional checks. Such an interface composition is still supported by
the base calculus, but the new cross-component handoffs, reads, key allocation,
and freshness obligations must be certified in the combined system.
Section~\ref{subsec:interface-composition} carries out exactly this step for the
leading interfering case: a shared append-only blackboard read by both
components together with cross-component handoffs and messages.
\end{remark}

\subsection{Interface Composition with Shared Reads}\label{subsec:interface-composition}

Separation forbids two features that realistic swarms need: cross-component
handoffs and a shared blackboard that both components can read. This subsection
weakens separation to the minimal discipline under which local certificates
still compose, and shows that the essential hypothesis is not the absence of
sharing but the \emph{append-only} status of shared memory together with
\emph{write-disjointness}. Cross-component reads are then admissible because an
append-only fact, once present, is stable under every later interleaving.

\begin{definition}[Interface composition with shared reads]\label{def:interface-composition}
Let \(\Gamma_k=(V_k,E_k)\), \(k=1,2\), with \(V_1\cap V_2=\emptyset\). An
\emph{interface composition} \(\Gamma_1\bowtie\Gamma_2\) consists of:
\begin{enumerate}
  \item the graph
  \((V_1\uplus V_2,\,E_1\uplus E_2\uplus E_{\times})\), where
  \(E_{\times}\subseteq(V_1\times V_2)\cup(V_2\times V_1)\) is a set of
  \emph{interface edges} carrying cross-component messages and handoffs, each
  with a configured payload policy;
  \item a single shared memory \(M\) over a tagged key space
  \(\Key=\Key_1\uplus\Key_2\) that is \emph{write-disjoint}: a step of
  component \(k\) may allocate only keys in \(\Key_k\), but any step of either
  component may \emph{read} any key in \(\operatorname{dom}(M)\);
  \item tagged, disjoint trace namespaces
  \(\Trace=\Trace_1\uplus\Trace_2\), and a shared certified-claim relation
  \(\kappa\) whose pairs \((a,k)\) may cite a key \(k\in\Key_1\cup\Key_2\)
  written by either component;
  \item per-component budgets and permissions
  \(\beta=\beta_1\sqcup\beta_2\) and \(\pi=\pi_1\uplus\pi_2\), and a product
  environment \(\Env=\Env_1\times\Env_2\).
\end{enumerate}
The combined ledger is
\(R_1\bowtie R_2=(\beta_1\sqcup\beta_2,\ \pi_1\uplus\pi_2,\
\iota_1\tau_1\cup\iota_2\tau_2,\ \kappa)\), with the shared memory \(M\) common
to both components rather than a product. A component step reads and writes the
shared memory subject to write-disjointness, and reads and writes only its own
budget, permission, trace, and environment factors. Interface steps along
\(E_{\times}\) are certified by the message and handoff schemas of the combined
system.
\end{definition}

The difference from separated composition
(Definition~\ref{def:swarm-composition}) is exactly the shared blackboard and
the interface edges: memory is one common append-only map that both components
read, and \(E_{\times}\) permits cross-component communication. Separation is the
special case \(E_{\times}=\varnothing\) with no cross-component read.

\begin{lemma}[Stability of append-only reads]\label{lem:appendonly-stability}
Let \(M_0\sqsubseteq M_1\sqsubseteq\cdots\sqsubseteq M_n\) be the shared-memory
states along any interleaving of write-disjoint component runs. If
\(k\in\operatorname{dom}(M_r)\) with \(M_r(k)=(v,p)\), then \(M_s(k)=(v,p)\) for
every \(s\ge r\). Consequently, every read premise and every certification
premise of the form \(k\in\operatorname{dom}(M)\) or \(M(k)=(v,p)\) discharged
at step \(r\) remains true at every later step.
\end{lemma}

\begin{proof}
Append-only extension gives \(M_s\sqsupseteq M_r\) for \(s\ge r\), so
\(k\in\operatorname{dom}(M_s)\); write-disjointness guarantees no component
overwrites a key allocated by the other, so no step can change the binding of
\(k\), whence \(M_s(k)=M_r(k)=(v,p)\). A premise asserting
\(k\in\operatorname{dom}(M)\) or \(M(k)=(v,p)\) therefore stays true once it
holds.
\end{proof}

\begin{theorem}[Compositionality under interface sharing]\label{thm:interface-compositionality}
Let each \((\Gamma_k,C_k,{\vdash}_{\Sigma_{\Gamma_k}})\) be locally sound, and
assume the interface conditions of
Definition~\ref{def:interface-composition}. Let a run of the composed system
interleave certified component steps with interface steps along \(E_{\times}\)
that are themselves certified by the combined message and handoff schemas. Then
the whole run is certified, and it satisfies
Theorem~\ref{thm:resource-preservation} for the combined ledger. In particular:
\begin{enumerate}
  \item each tool call decreases only its own component budget by \(\cost(t)\),
  and permissions remain component-local;
  \item the shared memory \(M\) remains a single append-only provenance map, and
  every cross-component read is stable in the sense of
  Lemma~\ref{lem:appendonly-stability};
  \item tagged trace identifiers remain globally fresh, and every certified
  claim---including a claim of one component citing a source key written by the
  other---stays supported in \(M\);
  \item every interface message or handoff respects an edge of \(E_{\times}\)
  and its payload policy.
\end{enumerate}
\end{theorem}

\begin{proof}
Consider one step of the interleaving. If it is a component step of component
\(k\), local soundness gives exact budget accounting, permission checking,
append-only writes into \(\Key_k\), a fresh tagged trace item, and claim support
against the current shared memory. Write-disjointness ensures its writes cannot
collide with keys of the other component, so \(M\sqsubseteq M'\) holds for the
common memory. If the step reads a key of the other component---either as a tool
input or as a certification source---Lemma~\ref{lem:appendonly-stability}
guarantees that the read fact is present and stable, so the local certificate
that discharged the read premise remains valid for the rest of the run. If the
step is an interface message or handoff along \(E_{\times}\), the combined
message or handoff schema discharges its graph and payload-policy premises and a
fresh tagged trace item; it changes no budget, memory, or claim, so it preserves
every invariant.

Cross-agent accessibility now genuinely crosses the components: by
Proposition~\ref{prop:graph-dependent-preorder}, an interface edge
\((i,j)\in E_{\times}\) puts \(\mathsf{agent}_i\preceq\mathsf{agent}_j\) in the
combined signature, so a rule focused on the receiving component may inspect the
sending component's reachable zone. This does not invalidate any certificate,
because the only cross-component resources a focused rule can use are reusable
facts in the shared, memory, and trace zones, all of which lie in \(U_{\Gamma}\)
and are monotone by Lemma~\ref{lem:appendonly-stability}; using such a fact
across the interface is a cut against a reusable formula, admissible by
Corollary~\ref{cor:cut-admissibility}, whereas the linear budget and obligation
tokens remain strictly component-local and are never accessed across the
interface. Induction on the length of the interleaving shows every step is
certified and every local-soundness clause holds for the combined ledger and
common memory. Applying Theorem~\ref{thm:resource-preservation} gives budget,
provenance, audit-freshness, and answer-support preservation for the whole run.
\end{proof}

\begin{corollary}[Separation as a special case]\label{cor:interface-generalises}
Separated composition (Definition~\ref{def:swarm-composition}) is the instance
of interface composition with \(E_{\times}=\varnothing\) and no cross-component
read. Hence Theorem~\ref{thm:interface-compositionality} strictly generalises
Theorem~\ref{thm:swarm-compositionality}: it additionally admits a shared
append-only blackboard read by both components and cross-component handoffs and
messages.
\end{corollary}

\begin{remark}[What interface sharing still excludes]
Two hypotheses remain essential and mark the boundary of the result.
Write-disjointness cannot be dropped: if either component could overwrite a key
allocated by the other, append-only monotonicity would fail and
Lemma~\ref{lem:appendonly-stability}, together with every cross-read
certificate, would become invalid. Likewise the shared claim relation
\(\kappa\) must stay monotone; a revocation event would reintroduce exactly the
interference the lemma rules out. Destructive update, deletion, key overwrite,
and claim revocation therefore still require combined re-verification and remain
outside the present theorem. What interface composition adds over separation is
precisely the realistic blackboard pattern---shared reads and cross-component
control transfer over append-only state---without a bialgebraic distributive
law.
\end{remark}

\section{Running Example}

The definitions above are parameterised by the agent set \(V\) and the interaction graph \(\Gamma\). The following three-agent system is a finite running instance, not the definition of a swarm. It is used as a concrete witness for the coalgebraic events, SELL-labelled resources, and checker artefacts.

Consider a research assistant swarm with three agents:
\[
  V=\{\mathsf{Planner},\mathsf{Retriever},\mathsf{Verifier}\}.
\]
The planner decomposes a research question, the retriever calls bibliographic tools, and the verifier checks whether claims are supported by sources. The interaction graph contains
\[
  \mathsf{Planner}\to\mathsf{Retriever},
  \qquad
  \mathsf{Retriever}\to\mathsf{Verifier},
  \qquad
  \mathsf{Verifier}\to\mathsf{Planner}.
\]
The finite checker instance assigns budgets
\[
  \beta_0(\mathsf{Planner})=3,
  \qquad
  \beta_0(\mathsf{Retriever})=5,
  \qquad
  \beta_0(\mathsf{Verifier})=2,
\]
permissions for \(\mathsf{BibliographicSearch}\) and \(\mathsf{ClaimCheck}\), empty shared memory \(M_0\), empty trace set \(\tau_0\), and empty certified-claim relation \(\kappa_0\). Thus
\[
  R_0=(\beta_0,\pi_0,\tau_0,\kappa_0),
  \qquad
  M_0=\varnothing,
  \qquad
  \tau_0=\varnothing,
  \qquad
  \kappa_0=\varnothing.
\]
The accepted trace \texttt{valid\_research\_swarm.json} then gives the certified execution in Table~\ref{tab:running-ledger}.

For example, the second row instantiates the certified transition relation. With \(q_2=\text{\texttt{tr-002}}\), \(k=\text{\texttt{src\_rutten2000}}\), and provenance \(p=(\mathsf{Retriever},\mathsf{BibliographicSearch},q_2)\), the raw call is admitted as
\[
  (y_1,R_1)
  \overset{\mathsf{call}(\mathsf{Retriever},\mathsf{BibliographicSearch},q_2),\,o_2}{\Longrightarrow}_{\cert}
  (y_2,R_2).
\]
The ledger and memory update are exactly
\[
  \beta_2(\mathsf{Retriever})=3,
  \qquad
  \tau_2=\tau_1\cup\{q_2\},
  \qquad
  M_2=M_1[k\mapsto(v,p)],
\]
while \(\pi\) and \(\kappa\) are unchanged. This concrete transition shows how the abstract judgement \(e:(y,R)\leadsto(y',R')\) is instantiated by the running ledger.

\begin{table}[htbp]
\centering
\scriptsize
\begin{tabular}{>{\raggedright\arraybackslash}p{0.08\linewidth}>{\raggedright\arraybackslash}p{0.17\linewidth}>{\raggedright\arraybackslash}p{0.29\linewidth}>{\raggedright\arraybackslash}p{0.30\linewidth}}
\toprule
Step & Event & Certificate premise & Ledger effect \\
\midrule
1 & handoff \(P\to R\) & graph edge, fresh trace & \(\tau\) gains \texttt{tr-001}; budgets unchanged \\
2 & search call by \(R\) & permission, budget \(5\ge2\), fresh trace & \(\beta(R)=3\); \(M\) gains \texttt{src\_rutten2000} with provenance \\
3 & message \(R\to V\) & graph edge, fresh trace & \(\tau\) gains \texttt{tr-003}; memory unchanged \\
4 & claim check by \(V\) & permission, budget \(2\ge1\), fresh trace & \(\beta(V)=1\); no new source key \\
5 & certify \((a,k)\) & authorised certifier, existing key \(k\) & \(\kappa\) gains \((a,k)\), with \(k=\texttt{src\_rutten2000}\) \\
6 & handoff \(V\to P\) & graph edge, fresh trace & certified payload returns to planner \\
7 & answer citing \(a\) & claim \(a\) already in \(\kappa\) & final answer is backed by \(k\); all trace identifiers are distinct \\
\bottomrule
\end{tabular}
\caption{Step-by-step ledger view of the running research-swarm example. Here \(P\), \(R\), and \(V\) abbreviate Planner, Retriever, and Verifier.}
\label{tab:running-ledger}
\end{table}

In the coalgebraic semantics, the seven rows form a path in the global raw swarm coalgebra. In the certified LTS, the same path exists only because every row has a matching SELL-labelled rule instance. The contrast with the negative traces is direct: if the planner attempts to call \(\mathsf{BibliographicSearch}\), the reusable permission premise is absent; if the verifier certifies a missing key, the shared-source premise is absent; if an answer cites an uncertified claim, the answer premise is absent. This is the intended division of labour: coalgebra describes the observable execution path, while SELL specifies which resources justify each step.

\section{Prototype Trace Checker}\label{sec:prototype}

To make the proposed discipline executable, we implemented a lightweight trace checker. The checker is located at \texttt{prototype/checker.py} in the accompanying repository~\cite{artifact2026}; it reads the finite swarm specification, validates the JSON traces, and writes the reproducible summaries used in this section. The generated artefacts are the JSON summary, the CSV summary, the LaTeX table input, a coverage summary, a coverage table, and a checksum manifest.

The prototype is not intended to be a complete SELL proof-search engine. Instead, it implements a decidable instance of the local-soundness conditions in Definition~\ref{def:local-soundness}. For each event, the implemented fragment checks whether:
\begin{itemize}
  \item the acting agent is known;
  \item handoffs and messages respect the interaction graph;
  \item tool calls are authorised by tool permissions;
  \item linear budgets are sufficient and are consumed exactly once;
  \item shared-memory writes are append-only and record provenance;
  \item final answers use only certified claims;
  \item tool calls, handoffs, messages, shared-memory updates, certifications, and answers have audit traces.
\end{itemize}

\begin{definition}[Abstract finite trace checker]\label{def:tracecheck}
For a finite swarm specification \(\mathcal{S}\) with agents, graph edges, tools, costs, permissions, certifiers, initial budgets, and auditable event kinds, define
\[
  \TraceCheck_{\mathcal{S}}
  : \mathsf{JSONTrace}_{\mathcal{S}}^{*}
  \longrightarrow
  (\{\Accept\}\times\mathsf{CheckState})
  + (\{\mathsf{Reject}\}\times\mathsf{Diag}^{*}).
\]
The function processes a finite event list \(T=e_1\cdots e_n\) from left to right. Its state is \((\beta,M,\kappa,\tau,D)\), consisting of the current budgets, shared memory, certified claim-source pairs, trace identifiers, and diagnostics. It starts from \((\beta_0,\varnothing,\varnothing,\varnothing,\varnothing)\). At step \(r\), it applies the corresponding one of the six certificate clauses: tool calls require a known agent, known tool, permission, sufficient budget, fresh audit item, exact budget decrement, and append-only writes; shared-memory updates require a known writer, fresh shared keys, admissible scope, provenance, and fresh audit item; handoffs and messages require known endpoints, a graph edge, policy admissibility, and fresh audit item; certifications require an authorised certifier, a pre-existing source key, and fresh audit item before adding \((a,k)\) to \(\kappa\); answers require every cited claim \(a\) to have some \((a,k)\in\kappa\) and a fresh audit item. Any failed premise appends a diagnostic to \(D\). The result is \(\Accept\) with the final state iff \(D\) is empty, and \(\mathsf{Reject}\) with diagnostics otherwise.
\end{definition}

\begin{algorithm}[t]
\caption{$\TraceCheck_{\mathcal{S}}(T)$: abstract finite trace checker}\label{alg:tracecheck}
\begin{algorithmic}[1]
\Require finite event sequence $T = e_1 \cdots e_n$; finite swarm specification $\mathcal{S}$
\Ensure $(\Accept, (\beta,M,\kappa,\tau))$ or $(\mathsf{Reject}, D)$
\State $(\beta, M, \kappa, \tau, D) \gets (\beta_0, \varnothing, \varnothing, \varnothing, \varnothing)$
\For{$r \gets 1$ \textbf{to} $n$}
  \If{$e_r = \mathsf{call}(i, t, q)$}
    \If{$i \notin \mathsf{Agents}(\mathcal{S})$ \textbf{or} $t \notin \mathsf{Tools}(\mathcal{S})$ \textbf{or} $(i,t)\notin\pi$ \textbf{or} $\beta(i) < \cost(t)$ \textbf{or} $q \in \tau$}
      \State $D \gets D \cup \{\mathsf{diag}(r,e_r)\}$
    \Else
      \State $\beta \gets \beta[i \mapsto \beta(i) - \cost(t)]$;\enspace $\tau \gets \tau \cup \{q\}$
    \EndIf
  \ElsIf{$e_r = \mathsf{writeShared}(i, k, v, p, q)$}
    \If{$i \notin \mathsf{Agents}(\mathcal{S})$ \textbf{or} $k \in \operatorname{dom}(M)$ \textbf{or} $\mathsf{scope}(k) \notin \mathsf{Scopes}(\mathcal{S})$ \textbf{or} $q \in \tau$}
      \State $D \gets D \cup \{\mathsf{diag}(r,e_r)\}$
    \Else
      \State $M \gets M[k \mapsto (v, p)]$;\enspace $\tau \gets \tau \cup \{q\}$
    \EndIf
  \ElsIf{$e_r = \mathsf{handoff}_{ij}(p,q)$ \textbf{or} $e_r = \mathsf{message}_{ij}(m,q)$}
    \If{$(i,j) \notin E_{\Gamma}$ \textbf{or} payload policy not satisfied \textbf{or} $q \in \tau$}
      \State $D \gets D \cup \{\mathsf{diag}(r,e_r)\}$
    \Else
      \State $\tau \gets \tau \cup \{q\}$
    \EndIf
  \ElsIf{$e_r = \mathsf{certify}(i, a, k, q)$}
    \If{$i \notin \mathsf{Certifiers}(\mathcal{S})$ \textbf{or} $k \notin \operatorname{dom}(M)$ \textbf{or} $q \in \tau$}
      \State $D \gets D \cup \{\mathsf{diag}(r,e_r)\}$
    \Else
      \State $\kappa \gets \kappa \cup \{(a, k)\}$;\enspace $\tau \gets \tau \cup \{q\}$
    \EndIf
  \ElsIf{$e_r = \mathsf{answer}(i, A, q)$}
    \If{$(\exists a \in A.\; \nexists k.\; (a,k) \in \kappa)$ \textbf{or} $q \in \tau$}
      \State $D \gets D \cup \{\mathsf{diag}(r,e_r)\}$
    \Else
      \State $\tau \gets \tau \cup \{q\}$
    \EndIf
  \EndIf
\EndFor
\If{$D = \varnothing$}
  \State \Return $(\Accept,\;(\beta,M,\kappa,\tau))$
\Else
  \State \Return $(\mathsf{Reject},\; D)$
\EndIf
\end{algorithmic}
\end{algorithm}

We write \(\TraceCheck_{\mathcal{S}}(T)=\Accept\) as shorthand for saying that the first component of the result is \(\Accept\).

The file \texttt{checker.py} (see~\cite{artifact2026}) is the reference implementation of \(\TraceCheck_{\mathcal{S}}\) for the JSON schema used in the artefact. The following propositions are stated about the abstract function \(\TraceCheck\); the Python program is evidence of reproducibility, not the mathematical object on which the proof depends. The file \texttt{rules.json} in the same repository is the concrete finite instance of the abstract specification~\(\mathcal{S}\); it is the object the checker reads at runtime, not a hypothesis in the mathematical sense.

\begin{proposition}[Soundness of accepted prototype traces]\label{prop:prototype-soundness}
For the finite swarm specification \(\mathcal{S}\) encoded by \texttt{rules.json} in~\cite{artifact2026}, every trace \(T\) such that \(\TraceCheck_{\mathcal{S}}(T)=\Accept\) determines a finite certified execution satisfying the local-soundness clauses of Definition~\ref{def:local-soundness}, with costs, permissions, graph edges, shared-memory keys, provenance records, certified claims, and trace items interpreted exactly as in \(\mathcal{S}\).
\end{proposition}

\begin{proof}
The function \(\TraceCheck_{\mathcal{S}}\) processes a trace from left to right while maintaining a ledger consisting of budgets, shared-memory keys with provenance records, certified claims, and trace identifiers. We show that, if no diagnostic is reported, every processed event satisfies the corresponding local-soundness clause.

First, for every auditable event, the trace clause checks that a trace identifier is present and has not previously appeared. Thus the transition from the current trace set \(\tau\) to the next trace set \(\tau'\) satisfies \(\tau\subseteq\tau'\), and for auditable events there is a fresh \(q\in\tau'\setminus\tau\).

Second, if the event is a tool call, the tool-call clause verifies that the acting agent is known, that the tool exists, that the agent belongs to the tool's allowed-agent list, and that the current budget is at least the tool cost. If these checks pass, \(\TraceCheck_{\mathcal{S}}\) subtracts exactly that cost from the agent's budget and leaves other agents' budgets unchanged. Hence the permission and budget clauses of Definition~\ref{def:local-soundness} hold with \(\rho_i(e)=0\) for the current prototype. Shared-memory writes are processed by the memory-write clause, which rejects missing keys, unsupported scopes, and attempts to overwrite an existing shared key; accepted writes therefore extend memory monotonically and attach the event, agent, tool when present, step, trace item, and value as provenance.

Third, if the event is a handoff or message, the corresponding clause verifies that both endpoint agents are known, the directed edge belongs to the interaction graph, and the payload satisfies the configured policy. Hence the graph and payload-policy conditions hold.

Fourth, if the event certifies a claim, the certification clause verifies that the certifying agent is authorised and that the cited source is already present in shared memory before adding the claim to the certified-claim set. Hence every newly certified claim is source-supported. If the event is a final answer, the answer clause verifies that each cited claim is already certified, so the answer condition holds.

The abstract checker records a diagnostic whenever one of these obligations fails. Therefore an accepted trace has no failed local-soundness obligation. Reading each accepted event as the corresponding certified transition yields the required finite certified execution.
\end{proof}

The following two tables are produced by running \texttt{checker.py} from the repository~\cite{artifact2026} against the full conformance suite. They are included verbatim as generated \LaTeX{} fragments; the generating command is given in Section~\ref{sec:reproducibility}.

% Generated by prototype/checker.py. Do not edit manually.
\begin{table}[htbp]
\centering
\scriptsize
\setlength{\tabcolsep}{3pt}
\begin{tabular}{@{}>{\raggedright\arraybackslash}p{0.24\linewidth}>{\raggedright\arraybackslash}p{0.10\linewidth}rr>{\raggedright\arraybackslash}p{0.48\linewidth}@{}}
\toprule
Trace & Status & Events & Cost & Diagnostic \\
\midrule
invalid\_\allowbreak{}bad\_\allowbreak{}memory\_\allowbreak{}scope & rejected & 1 & 0 & step 1: unsupported memory scope 'global' \\
invalid\_\allowbreak{}budget\_\allowbreak{}violation & rejected & 3 & 4 & step 3: insufficient budget for 'Retriever': required 2, available 1 \\
invalid\_\allowbreak{}certificate\_\allowbreak{}source & rejected & 1 & 0 & step 1: certificate source 'src\_\allowbreak{}missing' not found in shared memory \\
invalid\_\allowbreak{}duplicate\_\allowbreak{}trace & rejected & 2 & 0 & step 2: duplicate trace id 'tr-dup' \\
invalid\_\allowbreak{}handoff\_\allowbreak{}violation & rejected & 1 & 0 & step 1: handoff edge 'Planner -> Verifier' is not allowed \\
invalid\_\allowbreak{}message\_\allowbreak{}violation & rejected & 1 & 0 & step 1: message edge 'Planner -> Verifier' is not allowed \\
invalid\_\allowbreak{}missing\_\allowbreak{}trace & rejected & 1 & 2 & step 1: event 'tool\_\allowbreak{}call' lacks required audit trace \\
invalid\_\allowbreak{}shared\_\allowbreak{}memory\_\allowbreak{}overwrite & rejected & 2 & 4 & step 2: shared memory key 'src\_\allowbreak{}same' already exists \\
invalid\_\allowbreak{}tool\_\allowbreak{}permission & rejected & 1 & 2 & step 1: agent 'Planner' lacks permission for tool 'BibliographicSearch' \\
invalid\_\allowbreak{}unauthorized\_\allowbreak{}certifier & rejected & 2 & 2 & step 2: agent 'Planner' is not authorised to certify claims \\
invalid\_\allowbreak{}uncertified\_\allowbreak{}answer & rejected & 1 & 0 & step 1: answer uses uncertified claim 'unsupported\_\allowbreak{}claim' \\
invalid\_\allowbreak{}unknown\_\allowbreak{}agent & rejected & 1 & 0 & step 1: unknown source agent 'Ghost' \\
invalid\_\allowbreak{}unknown\_\allowbreak{}tool & rejected & 1 & 0 & step 1: unknown tool 'NonexistentTool' \\
valid\_\allowbreak{}claim\_\allowbreak{}reuse\_\allowbreak{}trace & accepted & 7 & 2 & OK \\
valid\_\allowbreak{}langgraph\_\allowbreak{}like\_\allowbreak{}trace & accepted & 7 & 3 & OK \\
valid\_\allowbreak{}memory\_\allowbreak{}update\_\allowbreak{}reuse\_\allowbreak{}trace & accepted & 6 & 0 & OK \\
valid\_\allowbreak{}research\_\allowbreak{}swarm & accepted & 7 & 3 & OK \\
valid\_\allowbreak{}two\_\allowbreak{}source\_\allowbreak{}research\_\allowbreak{}swarm & accepted & 9 & 5 & OK \\
\bottomrule
\end{tabular}
\caption{Results produced by the prototype trace checker. Accepted traces satisfy the graph, permission, budget, provenance, certification, and audit-trace constraints. Rejected traces illustrate violations of linear budget, certificate support, trace freshness, handoff topology, audit obligations, shared-memory append-only discipline, tool permission, and answer certification.}
\label{tab:prototype-results}
\end{table}

% Generated by prototype/checker.py. Do not edit manually.
\begin{table}[htbp]
\centering
\small
\begin{tabular}{>{\raggedright\arraybackslash}p{0.38\linewidth}>{\raggedright\arraybackslash}p{0.50\linewidth}}
\toprule
Metric & Value \\
\midrule
Total traces & 18 \\
Accepted traces & 5 \\
Rejected controls & 13 \\
Accepted event types & answer, certify, handoff, message, shared\_\allowbreak{}memory\_\allowbreak{}update, tool\_\allowbreak{}call \\
Accepted framework styles & langgraph-like \\
Accepted events & 36 \\
Accepted budget consumed & 13 \\
\bottomrule
\end{tabular}
\caption{Coverage summary for the finite conformance suite. The table reports suite size, accepted behaviours, exercised event kinds, and aggregate accepted-resource usage.}
\label{tab:prototype-coverage}
\end{table}

The accepted traces are minimal research-assistant conformance scenarios of the kind described above. They include a one-source trace, a two-source trace, a shared-memory-update trace, a claim-reuse trace showing that certified claims are reusable while audit identifiers remain fresh, and a LangGraph-like normalised trace carrying framework-style metadata. The rejected traces are negative controls covering budget overuse, unsupported certification, duplicate trace identifiers, illegal handoffs, payload-policy failure, missing audit evidence, shared-memory overwrite, unauthorised tool use, and answers that cite uncertified claims. Thus the artefact operationalises the formal claim behind Theorem~\ref{thm:resource-preservation}: by Proposition~\ref{prop:prototype-soundness}, accepted traces instantiate local soundness and therefore inherit graph, permission, budget, provenance, certification, and audit invariants. The suite is intentionally a finite conformance check, not evidence of empirical scalability.

Table~\ref{tab:obligation-checker-map} makes the correspondence explicit. Each formal obligation is discharged by one of the certificate schemas, implemented by a checker routine, and exercised by at least one positive or negative trace.

\begin{table}[htbp]
\centering
\scriptsize
\begin{tabular}{>{\raggedright\arraybackslash}p{0.20\linewidth}>{\raggedright\arraybackslash}p{0.19\linewidth}>{\raggedright\arraybackslash}p{0.20\linewidth}>{\raggedright\arraybackslash}p{0.29\linewidth}}
\toprule
Formal obligation & Certificate rule & Checker routine & Exercising traces \\
\midrule
Graph edge for handoffs/messages & \(\mathsf{handoff}\), \(\mathsf{message}\) & \texttt{check\_handoff}, \texttt{check\_message} & \texttt{valid\_research\_swarm}; \texttt{invalid\_handoff\_violation}, \texttt{invalid\_message\_violation} \\
Payload policy for handoffs/messages & \(\mathsf{handoff}\), \(\mathsf{message}\) & \texttt{payload\_allowed} & \texttt{valid\_research\_swarm}; \texttt{invalid\_payload\_policy} \\
Tool permission and known tool & \(\mathsf{call}\) & \texttt{check\_tool\_call} & \texttt{valid\_two\_source\_research\_swarm}; \texttt{invalid\_tool\_permission}, \texttt{invalid\_unknown\_tool} \\
Linear budget consumption & \(\mathsf{call}\) & \texttt{check\_tool\_call} & \texttt{valid\_research\_swarm}; \texttt{invalid\_budget\_violation} \\
Append-only source memory & \(\mathsf{writeShared}\) & \texttt{record\_writes} & \texttt{valid\_research\_swarm}; \texttt{invalid\_shared\_memory\_overwrite}, \texttt{invalid\_bad\_memory\_scope} \\
Fresh audit evidence & all auditable rules & \texttt{require\_trace} & \texttt{valid\_research\_swarm}; \texttt{invalid\_missing\_trace}, \texttt{invalid\_duplicate\_trace} \\
Claim-source certification & \(\mathsf{certify}\) & \texttt{check\_certify} & \texttt{valid\_two\_source\_research\_swarm}; \texttt{invalid\_certificate\_source}, \texttt{invalid\_unauthorized\_certifier} \\
Answer cites certified claims & \(\mathsf{answer}\) & \texttt{check\_answer} & \texttt{valid\_research\_swarm}; \texttt{invalid\_uncertified\_answer} \\
Known participating agents & all agent-indexed rules & \texttt{require\_known\_agent} & valid traces; \texttt{invalid\_unknown\_agent} \\
\bottomrule
\end{tabular}
\caption{Correspondence between formal obligations, certificate rules, checker routines, and conformance traces.}
\label{tab:obligation-checker-map}
\end{table}

For the implemented fragment, \(\TraceCheck\) is an executable recogniser for
admissible runs of the finite certificate calculus.

\begin{proposition}[Checker--calculus correspondence]\label{prop:checker-adequacy}
Fix the finite specification \(\mathcal{S}\) encoded by \texttt{rules.json}
in~\cite{artifact2026} and restrict events to the six certificate schemas above,
with no budget replenishment. A finite JSON trace \(T\) is accepted by
\(\TraceCheck_{\mathcal{S}}\) if and only if its events instantiate, in order,
admissible rules of the certificate calculus and produce the same ledger states.
Moreover, after erasing the external side-condition witnesses, the resource
transformers of that rule sequence form a focused derivation using the reusable
program clauses \(\mathcal{P}_{\mathcal{S}}\).
\end{proposition}

\begin{proof}
For the forward direction, read the accepted trace from left to right. Proposition~\ref{prop:prototype-soundness} shows that every accepted event satisfies the side conditions and resource premises of its corresponding certificate schema. Instantiate that schema at each step with the current ledger maintained by \(\TraceCheck_{\mathcal{S}}\). Its conclusion produces exactly the next checker state. The resource projection is one focused use of the corresponding reusable clause in \(\mathcal{P}_{\mathcal{S}}\); composing these uses yields the stated focused derivation.

For the converse direction, every admissible rule side condition is one of the predicates checked by the corresponding clause of Definition~\ref{def:tracecheck}. The rule conclusion and checker update agree on budgets, shared-memory extension, certified-claim insertion, and trace insertion. Induction over the rule sequence shows that no diagnostic is produced and that the maintained checker state coincides with the rule ledger. Hence the checker accepts the trace.
\end{proof}

\begin{proposition}[Complexity of the checker]\label{prop:checker-complexity}
Let \(s=|V|+|\Tool|+|E_{\Gamma}|+|\mathsf{Cert}|\) be the size of the
finite specification after parsing, and let \(N\) be the size of the trace
representation, including cited-claim and memory-write lists. With hash tables,
the acceptance pass of \(\TraceCheck_{\mathcal{S}}\) takes expected time
\(O(s+N)\) and space \(O(s+N)\), including preprocessing of the specification.
For fixed \(\mathcal{S}\), acceptance is therefore expected \(O(N)\). Producing
the reference implementation's deterministically sorted report can add
\(O(N\log N)\) time; this does not affect the acceptance decision.
\end{proposition}

\begin{proof}
Building agent, tool, edge, permission, and certifier hash tables takes expected
\(O(s)\) time and space. Each rule performs a bounded number of expected
constant-time membership, freshness, arithmetic, and update operations, plus
work linear in the event's explicit claims and writes. Summed across the trace,
that work is \(O(N)\). The maintained budgets and specification use \(O(s)\)
space, while memory, claims, traces, and diagnostics contain at most \(O(N)\)
data. Sorting report fields is separate from acceptance and gives the stated
optional \(O(N\log N)\) reporting bound.
\end{proof}

\begin{remark}
Proposition~\ref{prop:checker-adequacy} identifies exactly what the checker
decides: membership in the side-conditioned six-rule certificate language. The
bound in Proposition~\ref{prop:checker-complexity} follows from the local shape
of those rules. It is not a speed-up result for SELL proof search. For context,
Lincoln et al. show that MALL provability is \(\mathsf{PSPACE}\)-complete and
that full propositional linear logic is undecidable~\cite{lincoln1992}; those are
different decision problems, so no asymptotic comparison or optimality claim is
made here.
\end{remark}

The checker can be connected to existing agent frameworks by a trace-normalisation adapter rather than by changing the semantics. An AutoGen conversation log, a LangGraph execution state, or an OpenAI Agents SDK trace can be mapped into the prototype schema by extracting agent names, event kinds, directed handoffs or messages, tool invocations, memory writes, trace identifiers, and final-answer citations. This paper does not claim an empirical evaluation on those frameworks; the point is an interoperability path: framework traces become candidate raw executions, and the checker certifies the resource discipline when the required fields are present.

A stylised framework-normalisation step is shown in Table~\ref{tab:framework-normalisation}. It is not an additional experiment; it records the shape of the adapter assumed by the semantics. A LangGraph-style node transition or AutoGen message that invokes a tool is treated as a raw event first. Only after the permission, budget, memory, and audit fields are present can it become a certified event.

\begin{table}[htbp]
\centering
\scriptsize
\begin{tabular}{>{\raggedright\arraybackslash}p{0.25\linewidth}>{\raggedright\arraybackslash}p{0.35\linewidth}>{\raggedright\arraybackslash}p{0.28\linewidth}}
\toprule
Framework-level field & Normalised event field & Certificate obligation \\
\midrule
node or speaker name \texttt{Retriever} & \(\mathsf{agent}=\mathsf{Retriever}\) & known agent and local context \\
tool invocation \texttt{BibliographicSearch} & \(\mathsf{event}=\mathsf{tool\_call}\), \(\mathsf{tool}=\mathsf{BibliographicSearch}\) & reusable tool permission and linear budget \\
run/span identifier \texttt{tr-002} & \(\mathsf{trace}=q_2\) & fresh audit evidence \\
state write \texttt{src\_rutten2000} & shared write \((k,v,p)\) & append-only source memory with provenance \\
edge \texttt{Retriever -> Verifier} & message or handoff target \(\mathsf{Verifier}\) & graph admissibility \\
\bottomrule
\end{tabular}
\caption{Stylised normalisation of a framework trace step into the certificate schema.}
\label{tab:framework-normalisation}
\end{table}

\section{Assumptions, Reproducibility, and Limitations}\label{sec:reproducibility}

The preceding sections separate three layers: raw coalgebraic behaviour, SELL-labelled certification, and executable trace checking. This section records the assumptions under which the results hold, explains the role of the preservation theorem, and states the limits of the artefact. The purpose is to make the claims auditable rather than to enlarge them.

\subsection{Formal assumptions}

The deterministic core of the paper relies on the following assumptions.

\paragraph{Assumption 1: the base category is \(\Set\).}
States, observations, raw events, memories, ledgers, and traces are treated as sets, and the functors used in Propositions~\ref{prop:final-moore-semantics} and~\ref{prop:final-swarm-semantics} are endofunctors on \(\Set\). This conservative choice keeps the proof obligations explicit. Other bases, such as \(\mathbf{Rel}\) or a Kleisli category for probability, are meaningful extensions but are not used in the deterministic preservation theorem.

\paragraph{Assumption 2: the event functor is deterministic.}
The swarm functor is
\[
  S_{\Gamma}(X)=(\RawEv_{\Gamma}\times X)^{\Obs_{\Gamma}}.
\]
Thus a global state and a global observation determine one raw event and one next global state. The local agent interface
\[
  c_i:X_i\to(\Ev_i\times X_i)^{\Obs_i}
\]
is the local shape assembled into the global swarm coalgebra. A probabilistic variant would replace \(\RawEv_{\Gamma}\times X\) by \(\D(\RawEv_{\Gamma}\times X)\), but no theorem in this paper depends on that variant.

The genericity boundary is as follows. Existence of a behaviour map and
bisimulation invariance use only the chosen functor and its final coalgebra.
Likewise, for any endofunctor preserving the relevant inclusion, a subset whose
coalgebra structure factors through that inclusion is a subcoalgebra. The paper
instantiates this standard pattern for the deterministic Mealy functor because
its certificate rules inspect concrete events and ledgers. The distributional
shape above is an extension point, not a probabilistic theorem claimed here.

\paragraph{Assumption 3: SELL is used as a certification discipline.}
The paper fixes a concrete subexponential preorder \(\preceq_{\Gamma}\), a labelled context shape \(\Delta;\Theta\), and rule schemas for tool calls, memory writes, handoffs, messages, certification, and answers. These ingredients use SELL-labelled contexts to control which resources are reusable, local, consumable, or auditable. Table~\ref{tab:certificate-obligations} records the operational rule-to-obligation correspondence, and Lemma~\ref{lem:calculus-local-soundness} proves local soundness for the six schemas. The checker is not a full SELL proof-search engine. Instead, Definition~\ref{def:local-soundness} specifies the soundness obligations that any implementation of the rules must discharge: graph preservation, permission checking, budget accounting, shared-memory provenance, trace freshness, certification monotonicity, and answer safety. Proposition~\ref{prop:checker-adequacy} proves exact correspondence only for the finite implemented fragment. A full sequent-calculus implementation or proof-assistant mechanisation could replace the checker, provided these obligations are proved for each rule.

The relation to full SELL is deliberately delimited. The reusable program
\(\mathcal{P}_{\mathcal{S}}\) gives the resource transformer of every rule a
focused SELL interpretation, while graph membership, arithmetic, freshness,
and application policy are checked as external evidence. Proposition~\ref{prop:checker-adequacy}
establishes exact correspondence with these six side-conditioned schemas, not
soundness or completeness for arbitrary SELL formulae. An extension must add a
program clause, define its side conditions, and prove its local-soundness case.

\paragraph{Assumption 4: the preservation theorem is finite-run and safety-oriented.}
The behavioural results are stated over non-empty finite observation words, and Theorem~\ref{thm:resource-preservation} is stated for finite certified executions. The theorem does not prove liveness, optimality, convergence, or fairness. It does prove a prefix-closed safety property: by Corollary~\ref{cor:prefix-safety}, any long-running execution whose finite prefixes remain certified preserves the stated graph, permission, budget, provenance, certification, and audit invariants on every finite prefix.

\paragraph{Assumption 5: generic semantics and finitary checking are distinct.}
The final-coalgebra arguments are generic for the Set-functors used in the deterministic core. The prototype is finitary: it uses a finite set of agents, a finite tool catalogue, finite budgets, and finite JSON traces. At the semantic level, Propositions~\ref{prop:final-swarm-semantics} and~\ref{prop:behavioural-invariance} apply to arbitrary sets of observations and raw events. At the executable level, Proposition~\ref{prop:prototype-soundness} applies to the finite specification in \texttt{rules.json}~\cite{artifact2026}. The prototype illustrates and checks the theory; it does not replace the theory.

\subsection{Role of the preservation theorem}

The preservation theorem is intentionally not a deep theorem about all of SELL. Its role is to connect local certificates with global executions. Each certified transition must satisfy local obligations: the graph edge must exist, the tool must be permitted, the budget update must be exact, shared-memory writes must carry provenance, audit identifiers must be fresh, and final answers may cite only previously certified claims. Theorem~\ref{thm:resource-preservation} then proves that these local obligations compose over arbitrary finite runs. Consequently, a violation cannot be introduced merely by sequencing individually certified steps.

This local-to-global form is useful because raw coalgebraic behaviour and certified behaviour are intentionally different. The coalgebra describes what an implementation may emit; the certification layer identifies which emitted transitions are acceptable under the resource policy. The theorem states that the certified sub-behaviour has stable global invariants, even though the raw coalgebra may contain unauthorised, unaffordable, or unaudited transitions.

\begin{proposition}[Deterministic core under the stated assumptions]\label{prop:formal-closure}
Under Assumptions 1--5, every raw deterministic swarm coalgebra has a unique
observable behaviour map and raw-bisimilar states have identical observable
behaviours. Whenever \(X_{\cert}\) is closed under certified successors, it is
an \(S_{\Gamma}\)-subcoalgebra and certified bisimilarity is an equivalence on
\(X_{\cert}\) whose classes have identical raw behaviours. The agent-label
quotient of the SELL preorder is the reachability order of the graph's strongly
connected components. Certified paths form a prefix-closed safety sublanguage;
locally sound finite runs preserve resources and provenance; infinite runs are
prefix-safe; and locally sound swarms are closed under separated composition and
under interface composition with a shared append-only blackboard read by both
components. Modulo external side conditions, each certified event corresponds to
one focused phase of the certificate program.
Finally, accepted prototype traces instantiate the local-soundness conditions,
correspond exactly to admissible runs of the six-rule certificate calculus, and
are accepted in expected linear time in their representation size for fixed
\(\mathcal{S}\).
\end{proposition}

\begin{proof}
Each clause is exactly the corresponding previously proved proposition,
theorem, or corollary: Propositions~\ref{prop:final-swarm-semantics},
\ref{prop:behavioural-invariance}, \ref{prop:certified-subcoalgebra},
\ref{prop:graph-dependent-preorder}, \ref{prop:certified-safety-sublanguage},
\ref{prop:prototype-soundness}, \ref{prop:checker-adequacy}, and
\ref{prop:checker-complexity}; Theorems~\ref{thm:largest-certified-bisimulation},
\ref{thm:resource-preservation}, \ref{thm:swarm-compositionality},
\ref{thm:interface-compositionality}, and \ref{thm:focusing-adequacy}; and
Corollaries~\ref{cor:prefix-safety} and \ref{cor:composition-closure}.
\end{proof}

\subsection{Reproducibility artefact}

The artefact in \texttt{prototype/} is a conformance suite for the stated proof obligations, not a statistical benchmark of agent performance. It consists of \texttt{prototype/rules.json}, the trace files in \texttt{prototype/traces/}, the checker \texttt{prototype/checker.py}, the manifest verifier, and the generated outputs in \texttt{prototype/results/}; all files are publicly available at~\cite{artifact2026}. The checker uses the Python standard library only; Python 3.10 or newer is recommended. After cloning the repository, the complete regeneration and verification commands are:
\begin{verbatim}
git clone --branch jlamp-resubmission-2026-06 \
  https://github.com/cero1979/coalgebraic-agent-swarms
cd coalgebraic-agent-swarms
python3 prototype/checker.py \
  --rules prototype/rules.json \
  --traces prototype/traces/*.json \
  --out-dir prototype/results
python3 prototype/verify_manifest.py prototype/results/MANIFEST.json
\end{verbatim}
The expected conformance summary is:
\begin{itemize}
  \item 19 traces total;
  \item 5 accepted traces: \path{valid_research_swarm.json}, \path{valid_two_source_research_swarm.json}, \path{valid_memory_update_reuse_trace.json}, \path{valid_claim_reuse_trace.json}, and \path{valid_langgraph_like_trace.json};
  \item 14 rejected negative controls.
\end{itemize}
The first command writes \path{results.json}, \path{results.csv}, \path{results_table.tex}, \path{coverage_summary.json}, \path{coverage_table.tex}, and \path{MANIFEST.json} under \path{prototype/results/}; the second checks every recorded SHA-256 value. The outputs exclude runtime-version strings so regeneration is byte-stable across supported Python versions.
The generated LaTeX files \path{prototype/results/results_table.tex} and \path{prototype/results/coverage_table.tex} are included as Tables~\ref{tab:prototype-results} and~\ref{tab:prototype-coverage}.

The accepted traces exercise minimal certified research-assistant conformance scenarios, including separate shared-memory update, reusable certified-claim evidence, and LangGraph-like trace-normalisation metadata. The rejected traces exercise budget overuse, illegal handoff, payload-policy failure, missing audit evidence, duplicate trace identifiers, unsupported certificates, shared-memory overwrite, unauthorised tool use, and uncertified answers.

The artefact supports the formal development in four ways. First, it gives a concrete finite instance of the abstract ledger \(R=(\beta,\pi,\tau,\kappa)\). Second, it makes local-soundness failures observable as trace rejections. Third, accepted traces instantiate Proposition~\ref{prop:prototype-soundness}, and therefore inherit Theorem~\ref{thm:resource-preservation}. Fourth, Proposition~\ref{prop:checker-adequacy} identifies the checker as an executable recogniser for the implemented certificate fragment. This is evidence that the interface is executable, but it is not evidence of scalability, usability, or superiority over existing agent frameworks.

\subsection{Threats to validity and limitations}

The main limitations are as follows.
\begin{enumerate}
  \item The SELL layer is a side-conditioned certificate program rather than a complete proof-search implementation. Arithmetic, freshness, graph membership, and application policy are external evidence. The paper proves correspondence only for the implemented six-rule fragment; a mechanised development would need to encode the signature, program clauses, side-condition evidence, and local-soundness lemmas directly.
  \item Shared source memory is append-only in the deterministic core. This is appropriate for audit evidence, but it does not model destructive update, deletion, conflict resolution, or mutable heap ownership. Those features could be added using versioned keys, explicit revocation events, temporal provenance records, or a separation-logic-style ownership layer for mutable private state.
  \item The core theory is deterministic. Probabilistic policies, partial observability, reinforcement learning, MDP/POMDP objectives, and stochastic bisimulation require the distributional extension mentioned earlier but not developed here.
  \item The prototype uses a hand-written conformance suite supplemented by targeted negative mutations. It is sufficient to test the proof obligations implemented by the checker, but it is not a broad empirical evaluation of multi-agent systems.
  \item Claim certification is structural: it proves that an authorised certifier linked a claim to an existing provenance-carrying key. It does not prove the external truth of the claim or the semantic adequacy of the cited source.
  \item Interoperability with AutoGen, LangGraph, or the OpenAI Agents SDK is currently a trace-normalisation path rather than an implemented integration. Framework logs can be mapped to the schema when they expose agent names, event kinds, handoffs or messages, tool calls, memory writes, trace identifiers, and final-answer citations; evaluating such adapters is future work.
  \item The composition results cover separated components (Theorem~\ref{thm:swarm-compositionality}) and interface composition over a shared append-only blackboard read by both components (Theorem~\ref{thm:interface-compositionality}). They rely on write-disjointness and append-only monotonicity; shared mutable writes, key overwrite, and claim revocation break the stability of cross-component reads and still require combined re-verification.
\end{enumerate}

These limitations mark the boundary between the present contribution---a
deterministic semantics and resource-certification discipline with composition
only under separation---and the next stages: mechanised SELL proofs, verified
side-condition evidence, shared-state interface composition, richer memory
models, probabilistic coalgebras, and production-framework adapters.

\section{Discussion}

The results in this paper are contributions in the tradition of logical,
algebraic, and coalgebraic methods for the correctness and resource governance
of computing systems: the certified safety sublanguage result
(Proposition~\ref{prop:certified-safety-sublanguage}), the
resource-preservation theorem (Theorem~\ref{thm:resource-preservation}),
and the checker--calculus correspondence
(Proposition~\ref{prop:checker-adequacy}) are the primary logical-algebraic
results; the coalgebraic semantics and the SELL certification layer are the
formal machinery through which they are established. The following paragraphs
situate the contribution relative to closely related formalisms.

The proposed framework clarifies several distinctions.

First, an agent is not reduced to a cellular automaton. Cellular automata are specific local-update systems; agentic AI systems involve memory, tools, delegation, stochasticity, and external resources. Coalgebra is more general and better suited to the level of abstraction needed here.

Second, SELL is not proposed as a universal logic for all aspects of AI. It is proposed as a resource discipline. If the focus is mutable heap ownership, separation logic may be more appropriate. If the focus is epistemic change, dynamic epistemic logic is central. If the focus is static approximation of possible tool effects, a type-and-effect system may be the right abstraction. If the focus is optimal action under uncertainty, MDPs or POMDPs remain fundamental. If the focus is protocol fidelity and deadlock-freedom between communicating endpoints, session types and their linear-logic interpretation are the established tool~\cite{caires2010sessions,wadler2014sessions}; the present calculus instead leaves channel protocols implicit and controls quantitative budgets, reusable permissions, provenance, and audit obligations along certified executions. SELL is most useful when the central issue is controlled use of assumptions, budgets, permissions, provenance, and trace obligations.

Third, the formal contribution lies in the interface between coalgebra and resource logic. Coalgebra gives a semantics of possible behaviours; the side-conditioned SELL program recognises a certified sublanguage. The coupling is structural rather than bialgebraic: graph reachability determines the agent-zone preorder, while certificate rules use that signature to control context access. Framework guardrails can be understood as operational filters, while the certification calculus records why a transition is admissible and which linear or reusable resources justify it. Likewise, tracing changes role: it is not only a debugging record, but evidence tied to source keys, certified claims, and final answers.

Fourth, the framework explains why handoffs are more than routing edges. The graph condition says that control may move from one agent to another, but the certificate additionally records payload admissibility, trace evidence, and the ledger state available after transfer. Richer responsibility transfer could be modelled by adding explicit obligations to the handoff rule, for example duties to preserve a source key, revoke a stale claim, or emit a downstream verification trace.

Fifth, the conformance suite evaluates the implemented obligations, not agent performance or general SELL proof search. It exercises every implemented obligation with positive witnesses and targeted negative controls, and the checker--calculus correspondence shows that the code and six operational schemas accept the same traces. This methodology parallels certifying algorithms~\cite{mcconnell2011certifying}: evidence is checked independently of the system that produced it. Evaluating real AutoGen, LangGraph, or OpenAI SDK adapters is a separate future experiment.

\section{Conclusion}

This paper has applied logical, algebraic, and coalgebraic methods to the
correctness and resource governance of tool-augmented computing systems.
Individual agents are modelled as coalgebras from observations to events and
continuations. Swarms are modelled as global coalgebras indexed by an interaction
graph and equipped with shared memory, environment state, and a resource ledger.
SELL supplies a labelled proof-theoretic layer for distinguishing persistent
memory, shared context, local agent resources, tool permissions, consumable
budgets, and audit traces. LLM-based agent swarms provide the concrete motivating
instance, but the semantic framework applies to any computing system whose
component interactions are governed by explicit resource and certification policies.

The resulting framework offers a path toward formal reasoning about resource-governed computing systems: behavioural equivalence compares certifiable architectures, handoffs become graph-constrained certified transitions, and certified runs satisfy resource-preservation, provenance, and traceability properties. The assumptions and limitations above make explicit what is proved in the deterministic core: raw swarm coalgebras have canonical observable behaviours; certified bisimilarity on the closed certifiable subcoalgebra implies observational equality; certified paths form a prefix-closed safety sublanguage; locally sound runs preserve the stated invariants; and accepted prototype traces correspond to the six side-conditioned certificate schemas. Natural extensions are mechanisation, richer memory models, probabilistic coalgebras, mutation-based trace generation, and production-framework adapters.

The aim is not to replace existing agentic frameworks or empirical evaluations of them. It is to provide a compositional, resource-sensitive semantic layer on which framework traces can be interpreted, certified, and checked against explicit obligations for permissions, budgets, provenance, audit evidence, and answer support.

%% ---- Back matter (required by JLAMP / Elsevier) ----

\section*{CRediT authorship contribution statement}
\textbf{Carlos Ram\'irez Ovalle:} Conceptualization, Methodology,
Formal analysis, Software, Writing~-- original draft,
Writing~-- review and editing.

\printcredits

\section*{Declaration of competing interests}
The author declares that there are no competing interests.

\section*{Funding}
This research did not receive any specific grant from funding agencies
in the public, commercial, or not-for-profit sectors.

\section*{Data availability}
The prototype trace checker, conformance suite, and all generated
results are publicly available at
\url{https://github.com/cero1979/coalgebraic-agent-swarms}~\cite{artifact2026}.

\section*{Declaration of generative AI and AI-assisted technologies in the manuscript preparation process}
During revision, the author(s) used AI only for grammatical and idiomatic language adjustments. No AI tool was used to generate, verify, or modify the mathematical, computational, or scientific content of the article. The author(s) reviewed and edited all outputs and take full responsibility for the final manuscript.


%% ---- References ----

\begin{thebibliography}{99}
\raggedright

\bibitem{girard1987}
Jean-Yves Girard.
\newblock Linear logic.
\newblock \emph{Theoretical Computer Science}, 50(1):1--101, 1987.
\newblock DOI: \href{https://doi.org/10.1016/0304-3975(87)90045-4}{10.1016/0304-3975(87)90045-4}.

\bibitem{rutten2000}
Jan J. M. M. Rutten.
\newblock Universal coalgebra: A theory of systems.
\newblock \emph{Theoretical Computer Science}, 249(1):3--80, 2000.
\newblock DOI: \href{https://doi.org/10.1016/S0304-3975(00)00056-6}{10.1016/S0304-3975(00)00056-6}.

\bibitem{jacobs2016}
Bart Jacobs.
\newblock \emph{Introduction to Coalgebra: Towards Mathematics of States and Observation}.
\newblock Cambridge University Press, 2016.
\newblock DOI: \href{https://doi.org/10.1017/CBO9781316823187}{10.1017/CBO9781316823187}.

\bibitem{sokolova2011}
Ana Sokolova.
\newblock Probabilistic systems coalgebraically: A survey.
\newblock \emph{Theoretical Computer Science}, 412(38):5095--5110, 2011.
\newblock DOI: \href{https://doi.org/10.1016/j.tcs.2011.05.008}{10.1016/j.tcs.2011.05.008}.

\bibitem{pattinson2003}
Dirk Pattinson.
\newblock Coalgebraic modal logic: Soundness, completeness and decidability of local consequence.
\newblock \emph{Theoretical Computer Science}, 309(1--3):177--193, 2003.
\newblock DOI: \href{https://doi.org/10.1016/S0304-3975(03)00201-9}{10.1016/S0304-3975(03)00201-9}.

\bibitem{cirstea2011modal}
Corina Cîrstea, Alexander Kurz, Dirk Pattinson, Lutz Schröder, and Yde Venema.
\newblock Modal logics are coalgebraic.
\newblock \emph{The Computer Journal}, 54(1):31--41, 2011.
\newblock DOI: \href{https://doi.org/10.1093/comjnl/bxp004}{10.1093/comjnl/bxp004}.

\bibitem{hennessy1985}
Matthew Hennessy and Robin Milner.
\newblock Algebraic laws for nondeterminism and concurrency.
\newblock \emph{Journal of the ACM}, 32(1):137--161, 1985.
\newblock DOI: \href{https://doi.org/10.1145/2455.2460}{10.1145/2455.2460}.

\bibitem{sangiorgi2011}
Davide Sangiorgi.
\newblock \emph{Introduction to Bisimulation and Coinduction}.
\newblock Cambridge University Press, 2011.
\newblock DOI: \href{https://doi.org/10.1017/CBO9780511777110}{10.1017/CBO9780511777110}.

\bibitem{baier2008}
Christel Baier and Joost-Pieter Katoen.
\newblock \emph{Principles of Model Checking}.
\newblock MIT Press, 2008.

\bibitem{mcconnell2011certifying}
Ross M. McConnell, Kurt Mehlhorn, Stefan N{\"a}her, and Pascal Schweitzer.
\newblock Certifying algorithms.
\newblock \emph{Computer Science Review}, 5(2):119--161, 2011.
\newblock DOI: \href{https://doi.org/10.1016/j.cosrev.2010.09.009}{10.1016/j.cosrev.2010.09.009}.

\bibitem{necula1997pcc}
George C. Necula.
\newblock Proof-carrying code.
\newblock In \emph{Proceedings of the 24th ACM SIGPLAN-SIGACT Symposium on Principles of Programming Languages}, pages 106--119. ACM, 1997.
\newblock DOI: \href{https://doi.org/10.1145/263699.263712}{10.1145/263699.263712}.

\bibitem{cruzfilipe2017resolution}
Lu{\'i}s Cruz-Filipe, Jo{\~a}o Marques-Silva, and Peter Schneider-Kamp.
\newblock Efficient certified resolution proof checking.
\newblock In \emph{Tools and Algorithms for the Construction and Analysis of Systems}, LNCS 10205, pages 118--135. Springer, 2017.
\newblock DOI: \href{https://doi.org/10.1007/978-3-662-54577-5_7}{10.1007/978-3-662-54577-5\_7}.

\bibitem{leroy2009compcert}
Xavier Leroy.
\newblock Formal verification of a realistic compiler.
\newblock \emph{Communications of the ACM}, 52(7):107--115, 2009.
\newblock DOI: \href{https://doi.org/10.1145/1538788.1538814}{10.1145/1538788.1538814}.

\bibitem{baltag2003}
Alexandru Baltag.
\newblock A coalgebraic semantics for epistemic programs.
\newblock \emph{Electronic Notes in Theoretical Computer Science}, 82(1):17--38, 2003.
\newblock DOI: \href{https://doi.org/10.1016/S1571-0661(04)80630-3}{10.1016/S1571-0661(04)80630-3}.

\bibitem{cirstea2007}
Corina Cîrstea and Mehrnoosh Sadrzadeh.
\newblock Coalgebraic epistemic update without change of model.
\newblock In \emph{Algebra and Coalgebra in Computer Science, CALCO 2007}, Lecture Notes in Computer Science. Springer, 2007.

\bibitem{carreiro2013}
Facundo Carreiro, Daniel Gorín, and Lutz Schröder.
\newblock Coalgebraic announcement logics.
\newblock In \emph{Automata, Languages, and Programming: 40th International Colloquium, ICALP 2013}, pages 101--112. Springer, 2013.

\bibitem{hadzic2008}
Maja Hadzic and Elizabeth Chang.
\newblock Using coalgebra and coinduction to define ontology-based multi-agent systems.
\newblock \emph{International Journal of Metadata, Semantics and Ontologies}, 3(3):197--209, 2008.
\newblock DOI: \href{https://doi.org/10.1504/IJMSO.2008.023568}{10.1504/IJMSO.2008.023568}.

\bibitem{vanditmarsch2007}
Hans van Ditmarsch, Wiebe van der Hoek, and Barteld Kooi.
\newblock \emph{Dynamic Epistemic Logic}.
\newblock Springer, Synthese Library 337, 2007.
\newblock DOI: \href{https://doi.org/10.1007/978-1-4020-5839-4}{10.1007/978-1-4020-5839-4}.

\bibitem{feys2018}
Frank M. V. Feys, Helle Hvid Hansen, and Lawrence S. Moss.
\newblock Long-term values in Markov decision processes, (co)algebraically.
\newblock In \emph{Coalgebraic Methods in Computer Science, CMCS 2018}, LNCS 11202, pages 78--99. Springer, 2018.
\newblock DOI: \href{https://doi.org/10.1007/978-3-030-00389-0_6}{10.1007/978-3-030-00389-0\_6}.

\bibitem{cohen1990}
Philip R. Cohen and Hector J. Levesque.
\newblock Intention is choice with commitment.
\newblock \emph{Artificial Intelligence}, 42(2--3):213--261, 1990.
\newblock DOI: \href{https://doi.org/10.1016/0004-3702(90)90055-5}{10.1016/0004-3702(90)90055-5}.

\bibitem{shoham1993}
Yoav Shoham.
\newblock Agent-oriented programming.
\newblock \emph{Artificial Intelligence}, 60(1):51--92, 1993.
\newblock DOI: \href{https://doi.org/10.1016/0004-3702(93)90034-9}{10.1016/0004-3702(93)90034-9}.

\bibitem{wooldridge1995}
Michael Wooldridge and Nicholas R. Jennings.
\newblock Intelligent agents: Theory and practice.
\newblock \emph{The Knowledge Engineering Review}, 10(2):115--152, 1995.
\newblock DOI: \href{https://doi.org/10.1017/S0269888900008122}{10.1017/S0269888900008122}.

\bibitem{rao1995}
Anand S. Rao and Michael P. Georgeff.
\newblock BDI agents: From theory to practice.
\newblock In \emph{Proceedings of the First International Conference on Multi-Agent Systems (ICMAS 1995)}, pages 312--319, 1995.

\bibitem{wooldridge2009}
Michael Wooldridge.
\newblock \emph{An Introduction to MultiAgent Systems}.
\newblock Second edition, John Wiley \& Sons, 2009.

\bibitem{finin1994}
Tim Finin, Richard Fritzson, Don McKay, and Robin McEntire.
\newblock KQML as an agent communication language.
\newblock In \emph{Proceedings of the Third International Conference on Information and Knowledge Management}, pages 456--463. ACM, 1994.
\newblock DOI: \href{https://doi.org/10.1145/191246.191322}{10.1145/191246.191322}.

\bibitem{fipa2002}
Foundation for Intelligent Physical Agents.
\newblock FIPA ACL message structure specification.
\newblock Standard SC00061G, 2002.

\bibitem{nigam2009}
Vivek Nigam and Dale Miller.
\newblock Algorithmic specifications in linear logic with subexponentials.
\newblock In \emph{PPDP 2009: Proceedings of the 11th International ACM SIGPLAN Symposium on Principles and Practice of Declarative Programming}, pages 129--140. ACM, 2009.
\newblock DOI: \href{https://doi.org/10.1145/1599410.1599427}{10.1145/1599410.1599427}.

\bibitem{despeyroux2017}
Joëlle Despeyroux, Carlos Olarte, and Elaine Pimentel.
\newblock Hybrid and subexponential linear logics.
\newblock \emph{Electronic Notes in Theoretical Computer Science}, 332:95--111, 2017.
\newblock DOI: \href{https://doi.org/10.1016/j.entcs.2017.04.007}{10.1016/j.entcs.2017.04.007}.

\bibitem{andreoli1992}
Jean-Marc Andreoli.
\newblock Logic programming with focusing proofs in linear logic.
\newblock \emph{Journal of Logic and Computation}, 2(3):297--347, 1992.
\newblock DOI: \href{https://doi.org/10.1093/logcom/2.3.297}{10.1093/logcom/2.3.297}.

\bibitem{nigammiller2010}
Vivek Nigam and Dale Miller.
\newblock A framework for proof systems.
\newblock \emph{Journal of Automated Reasoning}, 45(2):157--188, 2010.
\newblock DOI: \href{https://doi.org/10.1007/s10817-010-9182-1}{10.1007/s10817-010-9182-1}.

\bibitem{reynolds2002}
John C. Reynolds.
\newblock Separation logic: A logic for shared mutable data structures.
\newblock In \emph{Proceedings of the 17th Annual IEEE Symposium on Logic in Computer Science (LICS 2002)}, pages 55--74. IEEE Computer Society, 2002.
\newblock DOI: \href{https://doi.org/10.1109/LICS.2002.1029817}{10.1109/LICS.2002.1029817}.

\bibitem{lucassen1988}
John M. Lucassen and David K. Gifford.
\newblock Polymorphic effect systems.
\newblock In \emph{POPL 1988: Proceedings of the 15th ACM SIGPLAN-SIGACT Symposium on Principles of Programming Languages}, pages 47--57. ACM, 1988.
\newblock DOI: \href{https://doi.org/10.1145/73560.73564}{10.1145/73560.73564}.

\bibitem{nielson1999}
Flemming Nielson, Hanne R. Nielson, and Chris Hankin.
\newblock \emph{Principles of Program Analysis}.
\newblock Springer, 1999.
\newblock DOI: \href{https://doi.org/10.1007/978-3-662-03811-6}{10.1007/978-3-662-03811-6}.

\bibitem{plotkin2003}
Gordon Plotkin and John Power.
\newblock Algebraic operations and generic effects.
\newblock \emph{Applied Categorical Structures}, 11(1):69--94, 2003.
\newblock DOI: \href{https://doi.org/10.1023/A:1023064908962}{10.1023/A:1023064908962}.

\bibitem{moreau2013prov}
Luc Moreau and Paolo Missier, editors.
\newblock PROV-DM: The PROV data model.
\newblock W3C Recommendation, 2013.

\bibitem{buneman2001}
Peter Buneman, Sanjeev Khanna, and Wang-Chiew Tan.
\newblock Why and where: A characterization of data provenance.
\newblock In \emph{Database Theory---ICDT 2001}, LNCS 1973, pages 316--330. Springer, 2001.
\newblock DOI: \href{https://doi.org/10.1007/3-540-44503-X_20}{10.1007/3-540-44503-X\_20}.

\bibitem{cheney2009}
James Cheney, Laura Chiticariu, and Wang-Chiew Tan.
\newblock Provenance in databases: Why, how, and where.
\newblock \emph{Foundations and Trends in Databases}, 1(4):379--474, 2009.
\newblock DOI: \href{https://doi.org/10.1561/1900000006}{10.1561/1900000006}.

\bibitem{puterman1994}
Martin L. Puterman.
\newblock \emph{Markov Decision Processes: Discrete Stochastic Dynamic Programming}.
\newblock John Wiley \& Sons, 1994.

\bibitem{kaelbling1998}
Leslie Pack Kaelbling, Michael L. Littman, and Anthony R. Cassandra.
\newblock Planning and acting in partially observable stochastic domains.
\newblock \emph{Artificial Intelligence}, 101(1--2):99--134, 1998.
\newblock DOI: \href{https://doi.org/10.1016/S0004-3702(98)00023-X}{10.1016/S0004-3702(98)00023-X}.

\bibitem{wu2023autogen}
Qingyun Wu et al.
\newblock AutoGen: Enabling next-gen LLM applications via multi-agent conversation framework.
\newblock \emph{arXiv preprint arXiv:2308.08155}, 2023.
\newblock DOI: \href{https://doi.org/10.48550/arXiv.2308.08155}{10.48550/arXiv.2308.08155}.

\bibitem{yao2023react}
Shunyu Yao et al.
\newblock ReAct: Synergizing reasoning and acting in language models.
\newblock In \emph{International Conference on Learning Representations}, 2023.

\bibitem{schick2023toolformer}
Timo Schick et al.
\newblock Toolformer: Language models can teach themselves to use tools.
\newblock In \emph{Advances in Neural Information Processing Systems}, 36, 2023.

\bibitem{li2024survey}
Xinyi Li, Sai Wang, Siqi Zeng, Yu Wu, and Yi Yang.
\newblock A survey on LLM-based multi-agent systems: Workflow, infrastructure, and challenges.
\newblock \emph{Vicinagearth}, 1:9, 2024.
\newblock DOI: \href{https://doi.org/10.1007/s44336-024-00009-2}{10.1007/s44336-024-00009-2}.

\bibitem{guo2024survey}
Taicheng Guo, Xiuying Chen, Yaqi Wang, Ruidi Chang, Shichao Pei, Nitesh V. Chawla, Olaf Wiest, and Xiangliang Zhang.
\newblock Large language model based multi-agents: A survey of progress and challenges.
\newblock \emph{arXiv preprint arXiv:2402.01680}, 2024.

\bibitem{openai2026agents}
OpenAI.
\newblock OpenAI Agents SDK documentation.
\newblock \url{https://openai.github.io/openai-agents-python/}, accessed May 12, 2026.

\bibitem{langgraph2026handoffs}
LangChain.
\newblock LangGraph multi-agent and handoff documentation.
\newblock \url{https://docs.langchain.com/oss/python/langchain/multi-agent/handoffs}, accessed May 12, 2026.

\bibitem{alpern1987safety}
Bowen Alpern and Fred B. Schneider.
\newblock Recognizing safety and liveness.
\newblock \emph{Distributed Computing}, 2(3):117--126, 1987.
\newblock DOI: \href{https://doi.org/10.1007/BF01782772}{10.1007/BF01782772}.

\bibitem{artifact2026}
Carlos Ram\'irez Ovalle.
\newblock Coalgebraic agent swarms: Prototype trace checker and conformance
  suite [software].
\newblock Pontificia Universidad Javeriana-Cali,
  \url{https://github.com/cero1979/coalgebraic-agent-swarms/tree/jlamp-resubmission-2026-06}, 2026.

\bibitem{lincoln1992}
Patrick Lincoln, John Mitchell, Andre Scedrov, and Natarajan Shankar.
\newblock Decision problems for propositional linear logic.
\newblock \emph{Annals of Pure and Applied Logic}, 56(1--3):239--311, 1992.
\newblock DOI: \href{https://doi.org/10.1016/0168-0072(92)90075-B}{10.1016/0168-0072(92)90075-B}.

\bibitem{milner1989}
Robin Milner.
\newblock \emph{Communication and Concurrency}.
\newblock Prentice Hall, 1989.

\bibitem{hoare1978csp}
C. A. R. Hoare.
\newblock Communicating sequential processes.
\newblock \emph{Communications of the ACM}, 21(8):666--677, 1978.
\newblock DOI: \href{https://doi.org/10.1145/359576.359585}{10.1145/359576.359585}.

\bibitem{milner1992picalculus}
Robin Milner, Joachim Parrow, and David Walker.
\newblock A calculus of mobile processes, I.
\newblock \emph{Information and Computation}, 100(1):1--40, 1992.
\newblock DOI: \href{https://doi.org/10.1016/0890-5401(92)90008-4}{10.1016/0890-5401(92)90008-4}.

\bibitem{caires2010sessions}
Lu\'is Caires and Frank Pfenning.
\newblock Session types as intuitionistic linear propositions.
\newblock In \emph{CONCUR 2010---Concurrency Theory}, LNCS 6269, pages 222--236. Springer, 2010.
\newblock DOI: \href{https://doi.org/10.1007/978-3-642-15375-4_16}{10.1007/978-3-642-15375-4\_16}.

\bibitem{wadler2014sessions}
Philip Wadler.
\newblock Propositions as sessions.
\newblock \emph{Journal of Functional Programming}, 24(2--3):384--418, 2014.
\newblock DOI: \href{https://doi.org/10.1017/S095679681400001X}{10.1017/S095679681400001X}.

\bibitem{honda1998}
Kohei Honda, Vasco T. Vasconcelos, and Makoto Kubo.
\newblock Language primitives and type discipline for structured communication-based programming.
\newblock In \emph{Programming Languages and Systems---ESOP 1998}, LNCS 1381, pages 122--138. Springer, 1998.
\newblock DOI: \href{https://doi.org/10.1007/BFb0053567}{10.1007/BFb0053567}.

\bibitem{huttel2016}
Hans H\"uttel, Ivan Lanese, Vasco T. Vasconcelos, Lu\'is Caires, Marco Carbone, Pierre-Malo Deni\'elou, Dimitris Mostrous, Luca Padovani, Ant\'onio Ravara, Emilio Tuosto, Hugo Torres Vieira, and Gianluigi Zavattaro.
\newblock Foundations of session types and behavioural contracts.
\newblock \emph{ACM Computing Surveys}, 49(1):3, 2016.
\newblock DOI: \href{https://doi.org/10.1145/2873052}{10.1145/2873052}.

\bibitem{das2018}
Ankush Das, Jan Hoffmann, and Frank Pfenning.
\newblock Work analysis with resource-aware session types.
\newblock In \emph{Proceedings of the 33rd Annual ACM/IEEE Symposium on Logic in Computer Science (LICS 2018)}, pages 305--314. ACM, 2018.
\newblock DOI: \href{https://doi.org/10.1145/3209108.3209146}{10.1145/3209108.3209146}.

\bibitem{turi1997}
Daniele Turi and Gordon Plotkin.
\newblock Towards a mathematical operational semantics.
\newblock In \emph{Proceedings of the 12th Annual IEEE Symposium on Logic in Computer Science (LICS 1997)}, pages 280--291. IEEE Computer Society, 1997.
\newblock DOI: \href{https://doi.org/10.1109/LICS.1997.614955}{10.1109/LICS.1997.614955}.

\bibitem{bonsangue2008}
Marcello M. Bonsangue, Jan Rutten, and Alexandra Silva.
\newblock Coalgebraic logic and synthesis of Mealy machines.
\newblock In \emph{Foundations of Software Science and Computational Structures, FoSSaCS 2008}, LNCS 4962, pages 231--245. Springer, 2008.
\newblock DOI: \href{https://doi.org/10.1007/978-3-540-78499-9_17}{10.1007/978-3-540-78499-9\_17}.

\end{thebibliography}

\end{document}
